\chapter{遥感与航拍}
\label{ch:20}

\paragraph*{本章阅读锚点}
本章的锚点是大视域遥感。遥感 3DGS 面对的是稀疏航拍、尺度跨度、地理坐标、正射视角和跨时间变化，核心难点不同于室内重建。


遥感与航拍图像的三维重建面临与地面场景截然不同的挑战：卫星采用 RPC（Rational Polynomial Camera，有理多项式相机）非线性相机模型而非标准针孔投影、影像跨越数百平方公里、多时相数据需要处理光照和季节变化。这些特殊性使得标准 3DGS 管线无法直接适用，需要在投影模型、光照建模和正则化等方面进行系统性适配。

本章系统梳理了~34~篇相关论文，按~4~个子方向展开讨论：卫星~3D~重建（第~\ref{sec:20_1}~节）、航空影像重建（第~\ref{sec:20_2}~节）、文化遗产与测绘（第~\ref{sec:20_3}~节）以及其他遥感方法（第~\ref{sec:20_4}~节）。

\begin{figure}[htbp]
\centering
\resizebox{\textwidth}{!}{\begin{tikzpicture}[
  route/.style={rectangle, rounded corners=5pt, draw=#1, fill=#1!10,
    text width=3.2cm, minimum height=0.9cm, align=center, font=\small},
  method/.style={rectangle, rounded corners=3pt, draw=#1, fill=#1!6,
    font=\scriptsize, text width=3.0cm, align=center, minimum height=0.6cm},
  arr/.style={-{Stealth[length=4pt]}, techgray!60},
  >=Stealth
]

  \node[route=techgray, text width=4cm, font=\small\bfseries] (root) at (22.1, 5.6) {遥感与航拍方法分类};
  \node[route=techblue] (b0) at (6.1, 3.1) {卫星3D重建};
  \node[method=techblue] (b0_l0) at (1.7, 0.6) {逐场景优化\\{\tiny EOGS, EOGS++, ShadowGS}};
  \node[method=techblue] (b0_l1) at (6.1, 0.6) {可泛化多场景\\{\tiny SkySplat, SatSurfGS, PSRGS}};
  \node[route=techgreen] (b1) at (17.5, 3.1) {航空影像重建};
  \node[method=techgreen] (b1_l0) at (10.7, 0.6) {航拍重建方法\\{\tiny ARSGaussian, AGS, Maikos et al.}};
  \node[method=techgreen] (b1_l1) at (15.2, 0.6) {文化遗产测绘\\{\tiny CULTURE3D, Gaussian Heritage, SkyLume}};
  \node[method=techgreen] (b1_l2) at (19.8, 0.6) {建筑精细重建\\{\tiny GBM, TOrtho-Gaussian, A-TDOM}};
  \node[route=techorange] (b2) at (24.4, 3.1) {农业与林业};
  \node[method=techorange] (b2_l0) at (24.4, 0.6) {作物表型分析\\{\tiny TreeDGS, Wheat3DGS, Li et al.}};
  \node[method=techorange] (b2_l1) at (28.8, 0.6) {生物量估计\\{\tiny Shen et al., Li et al.}};
  \node[route=techpurple] (b3) at (32.6, 3.1) {行星与水下};
  \node[method=techpurple] (b3_l0) at (32.6, 0.6) {行星遥感\\{\tiny Casadesus Vila et al., Moon-GS}};
  \node[method=techpurple] (b3_l1) at (36.2, 0.6) {水下遥感\\{\tiny Z-Splat, Qi et al.}};
  \draw[arr] (root) -- (b0);
  \draw[arr] (root) -- (b1);
  \draw[arr] (root) -- (b2);
  \draw[arr] (root) -- (b3);
  \draw[arr] (b0) -- (b0_l0);
  \draw[arr] (b0) -- (b0_l1);
  \draw[arr] (b1) -- (b1_l0);
  \draw[arr] (b1) -- (b1_l1);
  \draw[arr] (b1) -- (b1_l2);
  \draw[arr] (b2) -- (b2_l0);
  \draw[arr] (b2) -- (b2_l1);
  \draw[arr] (b3) -- (b3_l0);
  \draw[arr] (b3) -- (b3_l1);
\end{tikzpicture}}
\caption{遥感与航拍方法分类}
\label{fig:ch20_remote_sensing_taxonomy}
\end{figure}

%=====================================================================
\section{卫星3D重建}
\label{sec:20_1}

卫星~3D~重建是~3DGS~在遥感中最具挑战性的方向。卫星影像通常只有~5--20~个视角（远少于地面场景的数百甚至数千视角），采用有理多项式系数（RPC）相机模型（而非透视投影——RPC~将~3D~地理坐标映射到~2D~图像坐标，是一种非线性变换），且不同时相的影像间存在显著的光照和阴影差异。这些条件意味着标准~3DGS~的透视投影管线和~SfM（Structure from Motion，运动恢复结构）~初始化流程需要从根本上重新设计。

\subsection{逐场景优化方法}
\label{ss:20_1}

本节介绍逐场景优化的卫星~3DGS~方法。

\subsubsection{地球观测高斯Splatting}

本路线介绍将~3DGS~系统性适配到卫星~RPC~相机模型的逐场景优化方法。

将~3DGS~适配到卫星遥感场景需要解决投影模型转换、多时相光照不一致和稀疏视角正则化等核心问题。本节介绍的~EOGS~是首个系统性地将~3DGS~从地面透视场景扩展到卫星~RPC~相机模型的方法。

EOGS（地球观测高斯泼溅）（2024.12）~\cite{2412_13047}，全称~Earth Observation Gaussian Splatting，修改了标准~3DGS~的投影模型、光照建模和正则化机制——用仿射相机模型替代针孔透视投影、用高程差阴影映射替代无阴影渲染、用三重~Lagrangian~松弛约束替代无显式正则化——高斯基元的~splatting~渲染核心保持不变。本方法首次将~3DGS~系统性适配到遥感场景，解决了~3DGS~与卫星摄影测量的三个不兼容问题，在几分钟内达到最先进性能（最佳~NeRF~地球观测方法需一整天的优化）。

\paragraph{投影与坐标系。}标准~3DGS~假设针孔透视投影，即~3D~点通过内参矩阵线性映射到~2D~像素平面。然而卫星采用~RPC~相机模型，其投影是一个从~3D~地理坐标（经纬高）到~2D~图像坐标（行列号）的非线性有理多项式变换。整个变换链为：世界坐标$\to$UTM$\to$经纬高$\to$RPC$\to$图像行列$\to$NDC，无法直接套用~3DGS~的~EWA~splatting~公式。EOGS~的解决方案是将整条变换链近似为逐场景的仿射相机模型$A:\mathbb{R}^3 \to \mathbb{R}^2$，拟合后平均误差仅~0.012~像素。在此仿射近似下，3DGS~的投影公式获得本质简化：均值投影变为$\mu^A_k = A\mu_k + t \in \mathbb{R}^2$，协方差投影变为$\Sigma^A_k = A\Sigma_k A^\top \in \mathbb{R}^{2\times2}$，无需原版~3DGS~中~Jacobian~矩阵的局部一阶近似（因为近似已被转移到相机模型中）。

\paragraph{阴影映射。}标准~3DGS~不含任何光照建模——每个高斯的颜色由球谐系数直接编码，不区分直射光和阴影。然而卫星多时相影像中阴影广泛存在，忽略阴影会导致模型将阴影纹理"烘焙"到几何中（即半透明高斯通过调节不透明度模拟阴影，破坏几何精度）。EOGS~设计了兼容~Gaussian Splatting~的阴影映射方法。核心直觉是：如果一个~3D~点从太阳方向看和从卫星方向看"看到"的是同一海拔高度，则该点处于光照中；否则被遮挡，处于阴影中。

具体实现上，EOGS~首先定义高程渲染（用真实海拔替代颜色进行~3DGS~渲染）：$E^A(\mathbf{u}) = \sum_{k} \alpha^A_k(\mathbf{u}) \cdot E(\mathbf{x}_k)$，其中$E(\cdot)$将~3D~世界坐标映射为海拔（米）。然后构建与太阳方向对齐的"太阳相机"$S$，计算高程差$h^{A,S}(\mathbf{u}) = E^S(\mathbf{h}^{A,S}(\mathbf{u})) - E^A(\mathbf{u})$——将卫星视角的像素定位到~3D~空间后投影到太阳视角，比较两视角的高程渲染值。阴影系数为$s^{A,S}(\mathbf{u}) = \min\{\exp(-\rho \cdot h^{A,S}(\mathbf{u})),\, 1\}$，其中$\rho$为介质密度参数。逐相机环境光$\gamma^A$使阴影区域不至全黑，最终照明系数为$l^{A,S}(\mathbf{u}) = s^{A,S}(\mathbf{u}) + (1-s^{A,S}(\mathbf{u}))\gamma^A$。最终渲染输出为$I^{A,S}(\mathbf{u}) = l^{A,S}(\mathbf{u}) \cdot \sum_k c^A(\mathbf{c}_k)\,\alpha^A_k(\mathbf{u})$，其中$c^A(\cdot)$为逐相机仿射颜色校正（取代了原版~3DGS~的视角相关球谐颜色）。与标准~3DGS~的纯~$\alpha$-混合~渲染不同，此处渲染结果经过照明系数调制，阴影和光照效果被显式编码在管线中。这一设计避免了~EO-NeRF~等~NeRF~方法依赖的射线步进（ray marching，沿射线方向逐步推进检测遮挡）——射线步进在~Gaussian Splatting~中未被定义，因为高斯基元是离散的~3D~椭球体而非连续体密度场。

\paragraph{正则化。}3DGS~缺乏~NeRF~的隐式正则化——NeRF~的~MLP（Multi-Layer Perceptron，多层感知机）~在过参数化情况下仍能泛化，而~3DGS~的高斯基元几乎独立优化（它们初始化为均匀分布的小球），在稀疏视角下容易产生半透明伪影和冗余基元。EOGS~引入三种~Lagrangian~松弛约束：

（1）\textbf{$L^1$~不透明度稀疏性}——$L_o = \frac{1}{K}\sum_{k=1}^K |\alpha_k|$。受~LASSO~回归启发，通过$L^1$~范数鼓励高不透明度基元和低不透明度基元的两极分化，配合阈值剪枝（$\alpha < \alpha_{\min} = 0.0025$的基元被丢弃）减少基元数量，实测加速~$2\times$。

（2）\textbf{视角一致性}——对每个训练相机$A$随机扰动生成虚拟相机$B$（$B(\mathbf{x}) = A(\mathbf{x}) + 0.05\,E(\mathbf{x})[q_1,\, q_2]^\top$，$q_1,q_2 \sim \text{truncated-}\mathcal{N}(0,1)$），要求两者在高程差$h^{A,B}(\mathbf{u}) < h_{\max} = 30$~m~的区域中渲染一致：颜色一致性损失$L_{cc} = \sum_{\mathbf{u}} M^{A,B}(\mathbf{u})\,|I^A(\mathbf{u}) - I^B(\mathbf{h}^{A,B}(\mathbf{u}))|$和高程一致性损失$L_{ac} = \sum_{\mathbf{u}} M^{A,B}(\mathbf{u})\,|E^A(\mathbf{u}) - E^B(\mathbf{h}^{A,B}(\mathbf{u}))|$，无需预训练模型或深度先验。

（3）\textbf{不透明度熵惩罚}——$L_s = \sum_{\mathbf{u}} H(s^{A,S}(\mathbf{u}))$，其中$H(x) = -(x\log_2 x + (1-x)\log_2(1-x))$。动机是阴影图应只包含~0（全阴影）或~1（全光照），半透明阴影意味着半透明物体——而建筑物应投射完全不透明的阴影。熵惩罚迫使阴影图趋向二值化，间接迫使高斯基元变为完全不透明或完全透明。

总损失函数为$\sum_{i=1}^N \mathcal{L}(I_i, \hat{I}_i) + 0.1\,L_o + 0.1\,L_{cc} + 0.01\,L_{ac} + 0.01\,L_s$，其中四个~Lagrangian~系数在单场景上调定后取最近的~10~的幂次，即可泛化到所有场景。实现上，EOGS~还禁用了原版~3DGS~的视角相关球谐颜色（改为逐相机仿射颜色校正），所有高斯初始化为白色且不透明度仅~1\%，基元密度为~0.13~个/m$^3$，总共仅~5000~次迭代（阴影和正则化从第~1000~次迭代起启用）。

EOGS~虽然将~3DGS~系统性适配到了遥感场景，但其管线仍然依赖两个外部预处理步骤：基于~SIFT~的光束法平差（BA，Bundle Adjustment）修正~RPC~姿态的地理定位误差（无~BA~时~MAE~从~1.33~m~恶化至~1.93~m），以及全色锐化（pan-sharpening）将低分辨率多光谱影像融合为~RGB~图像以适配~3~通道渲染。这两步增加了工程复杂度且可能引入对齐误差。EOGS++（2025.11）~\cite{2511_16542} 的核心目标是消除这些外部预处理依赖，同时在训练和后处理环节进一步改进，共涉及四方面改动。

第一，将光束法平差内嵌到训练循环中。EOGS++~用~RAFT-Small~光流估计器计算渲染图像与训练图像之间的逐像素位移场，取其均值得到常数偏移~$\bar{\Delta} \in \mathbb{R}^2$~并对渲染图像做双线性重采样~$\tilde{I}^A(\mathbf{u}) = I^A(\mathbf{u}+\bar{\Delta})$；光流梯度被~stop-gradient~截断，确保反向传播不流经光流模块，从而隐式修正相机姿态。该方法略逊于外部~BA（平均~MAE~差距仅~0.03~m），但完全消除了外部预处理依赖。

第二，用"3-PAN"策略取代全色锐化预处理。直接将单通道全色图像复制为三通道~RGB~表示，与~EOGS~的逐相机仿射颜色变换$c^A(\cdot)$直接兼容（无需修改颜色校正模块），省略了多光谱融合步骤。实验表明~3-PAN~的重建精度与使用~Brovey~全色锐化的~EOGS~完全一致，说明对~DSM~估计任务而言多光谱信息增益有限。

第三，重新引入原版~3DGS~的不透明度重置机制（每~3000~次迭代将所有高斯不透明度截断为~$\min(\alpha_k,\, 0.05)$）以消除浮游高斯，同时将最大迭代次数从~EOGS~的~5000~提升至~40000。然而不透明度重置导致光度损失呈现非单调行为——先下降后上升（推测原因是~EOGS~的正则化项在训练后期主导了优化方向），因此~EOGS++~设计了基于光度损失极小值的早停机制：训练在光度损失首次回升时自动终止，实际收敛发生在~10000$\sim$20000~次迭代。

第四，用~TSDF（Truncated Signed Distance Function，截断符号距离函数）~后处理替代~EOGS~的虚拟天底相机渲染。EOGS~原用一个未参与训练的虚拟天底相机生成~DSM（Digital Surface Model，数字表面模型），当训练视角偏离天底构型时易产生伪影。EOGS++~改为用优化后的高斯基元从所有训练视角渲染深度图，按视角方向与估计表面法向量之间夹角加权融合为~TSDF~体素网格（截断边距固定），经~Marching Cubes~提取网格后去除孤立体素并自底向上填充孔洞，最终在地面位置采样高度得到~DSM。

在~IARPA 2016~和~DFC2019~两个公开遥感数据集上，EOGS++~将建筑物平均~MAE~从~EOGS~的~1.33~m~降至~1.19~m，训练时间约~25~分钟，远快于~NeRF~方法（EO-NeRF~需~900~分钟），在重建精度和计算效率上均取得了显著改进。

ShadowGS（2026.01）~\cite{2601_00939} 修改了~3DGS~的阴影计算、光照表示和几何建模三个环节——用硬件加速光线步进（ray-marching）替代~EOGS~的太阳相机阴影映射，用反照率--光照双~SH~系数的物理渲染方程替代~EOGS~的逐相机环境光补偿，用~RaDe-GS~射线--高斯交集替代标准~splatting~投影。~3DGS~的~$\alpha$-混合~渲染框架和高斯基元参数化保持不变。

\paragraph{几何表示：射线--高斯交集} 在几何表示上，ShadowGS~遵循~RaDe-GS~的显式射线--高斯交集策略，而非标准~3DGS~的纯~splatting~投影。直觉上，标准~3DGS~只关心每个高斯在像素上的~2D~投影贡献，不区分该像素看到的到底是高斯的"顶部"还是"侧面"；而射线--高斯交集策略相当于对每条视线与每个高斯椭球体的交线做精确求解，从而为每个像素获得精确的深度值和表面法向量。具体而言，对于~3D~高斯$G$，其~2D~投影中心为$(u_c, v_c)$，像素$(u, v)$处射线与高斯的交集点定义为一维高斯分布的峰值位置，对应的深度贡献为
\begin{equation}
d = z_c + \frac{\sigma_{ct_c} \, m^\top \Sigma'^{-1} \begin{pmatrix} u_c - u \\ v_c - v \end{pmatrix}}{\sigma_{ct_c} \cdot v^\top \Sigma'^{-1} v}
\end{equation}

其中$z_c$和$t_c$分别为高斯中心的深度和高斯中心到相机中心的距离，$m$为从$2\times 2$投影协方差矩阵相关构造中得到的$1\times 2$向量。该交集点处的表面法向量定义为$n = J^\top (-m - e_1)^\top$（$J$为投影雅可比矩阵），归一化深度图$D$和法向图$N$通过标准~$\alpha$-混合~渲染。为增强几何细节，ShadowGS~进一步引入深度--法向一致性损失$L_n = (1 - N^\top \hat{N})$，其中$\hat{N}$为从渲染深度图$D$通过有限差分推导的表面法向量——若法向直接来自几何，则$N$与$\hat{N}$应一致，此损失驱使高斯基元更好地贴合真实表面。

\paragraph{阴影计算：硬件加速光线步进} EOGS~的阴影映射本质上是一种间接近似：通过在太阳方向构造虚拟"太阳相机"比较高程渲染差异来推断阴影状态，在复杂遮挡（如多层建筑、密集城区）下容易遗漏或误判。ShadowGS~采用硬件加速光线步进（ray-marching）策略替代阴影映射，实现像素级精确的阴影计算。核心思想直观理解：从每个高斯中心沿太阳方向"逆光前行"，沿途检测是否被其他高斯遮挡——如果前方有足够多的高斯阻挡阳光，该点就处于阴影中。

在实现上，所有高斯基元被组织为自适应缩放的有界体层次（BVH），具体采用拉伸二十面体（stretched icosahedron）包围盒，其顶点$a$经过如下变换以适应每个高斯的有效范围：

\begin{equation}
a' = a \cdot 2(o / o_{\min})^{1/2} \cdot S R^\top + \mu
\end{equation}

其中$o$为高斯不透明度，$o_{\min} = 0.001$为透明度阈值（确保包围体完全覆盖高斯有效区域），$S$和$R$分别为缩放和旋转矩阵，$\mu$为高斯中心位置。这一缩放策略使不透明度高（对光线阻挡强）的高斯拥有更大的包围盒，从而被光线步进更早检测到。

对于每个高斯中心$\mu_0$（作为射线起点），沿太阳方向$r$以固定步长投射光线。光线与第$i$个高斯的交集点$\tau$为一维高斯分布的峰值位置：

\begin{equation}
\tau = \frac{(\mu_0 - \mu_i)^\top \Sigma_i^{-1} r}{r^\top \Sigma_i^{-1} r} \cdot r
\end{equation}

该点处的高斯响应值（即对光线的遮挡强度）为$\beta_i = o_i \cdot G^{1D}(\tau)$，其中$G^{1D}(\tau)$为交集一维高斯分布在$\tau$处的概率密度。太阳可见度$S_{\text{sun}}$为所有遮挡物透射率的连乘积：

\begin{equation}
S_{\text{sun}} = \prod_{i=1}^{k}(1 - \beta_i)
\end{equation}

直觉上，$S_{\text{sun}} = 1$表示阳光直达该点（无遮挡），$S_{\text{sun}} \to 0$表示完全处于阴影中。每个遮挡高斯按其不透明度和空间位置贡献遮蔽量，遮挡越密集，可见度越低。

\paragraph{物理渲染方程：光照与外观解耦} EOGS~使用逐相机环境光$\gamma^A$补偿阴影区域，本质上是一个标量常数，无法区分阴影中的不同光照来源（如天空散射、附近建筑反射）。ShadowGS~设计了完整的物理渲染方程，显式解耦直射阳光、天空散射和近表面反射三种光照分量。

具体而言，每个高斯配备两组球谐（SH）系数：第一组编码反照率（albedo）$f$——物体表面的固有颜色，与光照无关；第二组编码近表面反射辐射$l_n$——邻近表面反射到该点的间接光照。此外，所有高斯共享一组全局天空光球谐$l_{\text{sky}}$（仅启用低阶项，限制学习到低频的空间均匀天空辐射）。反照率$F$和各光照分量$(S, L_{\text{sky}}, L_n)$分别通过~$\alpha$-混合~渲染为~2D~图。每个表面点的总入射辐射建模为
\begin{equation}
L_{\text{total}} = S_{\text{sun}} \cdot S + (1 - S_{\text{sun}})(L_{\text{sky}} + L_n)
\end{equation}

直觉理解：阳光直射区域由直射阳光$S$主导（$S_{\text{sun}} \approx 1$），阴影区域则由天空散射$L_{\text{sky}}$和近表面反射$L_n$共同照亮（$S_{\text{sun}} \approx 0$）。最终渲染颜色为反照率与总入射辐射的逐元素乘积：

\begin{equation}
C = F \cdot L_{\text{total}}
\end{equation}

这一公式的物理含义是：表面颜色 = 固有颜色 $\times$ 入射光照——解耦后的反照率在多时相影像间保持一致，光照变化则由$S_{\text{sun}}$和$l_{\text{sky}}$等参数吸收。

\paragraph{阴影一致性约束} 3DGS~的离散高斯表示比~NeRF~的连续神经表示更加不规则，常导致优化约束不足。ShadowGS~利用卫星成像的独特几何特性引入阴影一致性约束。直觉上，卫星和太阳都距地面极远，当卫星视角与太阳方向共线时，物体投射的阴影恰好被物体自身遮挡——即此时图像中应看不到任何阴影。形式化地，给定虚拟相机视点和与之共线的太阳方向，渲染阴影图$S_v$，阴影一致性损失定义为
\begin{equation}
L_{S_1} = \| S_v - 1 \|_1
\end{equation}

即驱使共线条件下的阴影图趋向全亮（$S_v = 1$）。实践中采用两种配置实现：（1）固定相机视点，将太阳方向对齐到视线方向；（2）同时调整太阳和相机方向使其垂直于场景表面。此约束鼓励表面高斯达到更高的不透明度并更好地贴合几何表面。

ShadowGS~还引入了阴影二值化损失$L_{S_2} = -[S \cdot \log_2(S) + (1-S) \cdot \log_2(1-S)]$（与~EOGS~一致），鼓励投影阴影趋向完全透明或完全不透明的二值状态。

\paragraph{稀疏视角下的阴影先验} 稀疏输入视角（如仅~5--10~张卫星影像）对~3DGS~和~NeRF~均构成严峻挑战，常导致过拟合和性能退化。现有方法多依赖深度先验或扩散模型，前者在纹理缺乏区域失效，后者可能引入错误的伪视角。ShadowGS~的创新在于利用几何相关的阴影信息作为鲁棒监督信号——阴影的形状和位置严格由场景几何决定，因此可作为几何优化的间接约束。具体地，集成~FDRNet~自监督阴影检测网络（在卫星影像中假阴性率低），通过二元交叉熵损失最小化渲染阴影图$S$与~FDRNet~提取的阴影掩码$\Psi$之间的差异：

\begin{equation}
L_{S_3} = -[S \cdot \log_2(\Psi) + (1 - S) \cdot \log_2(1 - \Psi)]
\end{equation}

为处理植被区域的虚假阴影检测，采用归一化植被指数（NDVI，多光谱数据）或差值增强植被指数（DEVI，仅~RGB~数据）排除植被区域的阴影监督；并在致密化结束后（第~3000~次迭代后）停止阴影先验损失，以防止残余假阳性干扰优化。

\paragraph{训练策略与完整损失} ShadowGS~的总损失函数为
\begin{equation}
\mathcal{L} = \lambda_c \mathcal{L}_c + \lambda_n \mathcal{L}_n + \lambda_{S_1} \mathcal{L}_{S_1} + \lambda_{S_2} \mathcal{L}_{S_2} + \lambda_{S_3} \mathcal{L}_{S_3}
\end{equation}

其中$\mathcal{L}_c = (1-\lambda_{\text{ssim}})\mathcal{L}_{\text{color}} + \lambda_{\text{ssim}}\mathcal{L}_{\text{D-SSIM}}$为标准光度重建损失。在~DFC2019~数据集上，各损失权重设为$\lambda_c = 10$、$\lambda_n = 0.5$、$\lambda_{S_1} = 0.2$、$\lambda_{S_2} = 0.3$、$\lambda_{S_3} = 1$；在~IARPA~数据集上$\lambda_n$降至~0.1。训练策略相比标准~3DGS~有关键调整：全程禁用不透明度重置（避免破坏已解耦的光照参数）、所有损失项从第~1~次迭代即激活、总训练~5000~次迭代（致密化在第~3000~次后停止）、致密化阶段每~300~次迭代执行一次（标准~3DGS~为每~100~次）——降低致密化频率使高斯基元在被进一步分裂前获得更充分的优化，减轻冗余高斯增殖和计算开销。阴影先验损失$L_{S_3}$仅在稀疏视角条件下激活。

与~SkySplat~同样采用可泛化前馈框架和~RPC~代价体构建高度图，SatSurfGS（2026.05）~\cite{2605_07181} 的核心差异在于两方面：表示形式和可靠性建模。首先，SatSurfGS~用~2D Gaussian Surfel~替代~3D~高斯基元——卫星自上而下观测地表，地形可由~2.5D~高度图充分描述，2D Surfel~的平面假设天然契合这一先验，避免了~3D~椭球体在深度方向的不确定性；其次，SatSurfGS~观察到稀疏卫星视角下多视角几何约束的可靠性存在空间异质性（弱纹理、辐射不一致区域约束弱，纹理丰富区域约束强），由此提出以代价体置信度作为统一可靠性描述符，贯穿特征融合、参数估计和训练监督三个层面。SatSurfGS~采用三阶段（$s=1,2,3$）由粗到细级联设计，代价体构建沿用~RPC~反投影 + 多视角特征方差的范式（与~SkySplat~一致），但新增三个置信度驱动模块：

（1）\textbf{代价体置信度}——对代价体沿高度维度做~softmax~归一化后，取概率分布的最大响应值$C_s = \max_k P_s^k$作为局部几何可靠性指标。相比光度残差（受光照和阴影干扰）、图像梯度（只描述外观丰富性，不反映跨视角一致性）等替代方案，最大概率响应直接衡量多视角匹配是否有确定性的唯一高度解——分布尖锐（高置信）意味着多视角一致支持某个高度假设，分布平坦（低置信）则表示匹配模糊。实验分析表明，$C_s$与高度误差强相关：随着平均置信度升高，MAE~和~RMSE~持续下降。

（2）\textbf{置信度感知的单目-多视角特征融合（CMMF）}——以多视角匹配特征$V_{s,\text{mvs}}^j$为主干，通过置信度调制的注意力偏置动态引入单目先验特征$V_{s,\text{mono}}^j$的残差修正。具体地，对双分支分别计算注意力对数$\lambda_{s,\text{mvs}}^j = Q_s^{j\top}K_{s,\text{mvs}}^j/\sqrt{d} + \phi(\bar{C}_{s-1})$和$\lambda_{s,\text{mono}}^j = Q_s^{j\top}K_{s,\text{mono}}^j/\sqrt{d} - \phi(\bar{C}_{s-1})$，其中$\phi(\cdot)$为置信度生成的偏置项（含缩放因子和阈值），$\bar{C}_{s-1}$为前一阶段上采样后的置信度图。高置信区域增强~MVS~分支、抑制单目分支，低置信区域则反转。融合采用残差路由而非线性平均：$\hat{F}_s^j = V_{s,\text{mvs}}^j + W_{s,\text{mono}}^j \odot (V_{s,\text{mono}}^j - V_{s,\text{mvs}}^j)$，单目分支以修正残差而非直接替换的方式引入；且单目分支的值投影层初始化为零，训练早期模块近似恒等映射，防止单目先验干扰多视角几何。

（3）\textbf{跨阶段自洽性残差引导（CSRG）}——计算前一阶段~2DGS~渲染高度图$\hat{H}_{s-1}$与当前阶段~MVS~估计高度图$H_s$之间的鲁棒归一化残差。直接差值不稳定（不同阶段的高度分布可能存在尺度偏移），因此先做鲁棒归一化$\tilde{H} = (H - \text{Median}(H))/(\text{MAD}(H) + \epsilon)$（MAD~为中位绝对偏差），再计算残差$\Delta H_s = \tilde{H}_s - \text{up}(\tilde{H}_{s-1})$。将方向信息$\Delta H_s$、幅度信息$|\Delta H_s|$和当前阶段置信度$C_s$拼接为三通道引导图$R_s$，通过注意力控制的残差注入当前阶段的高斯参数预测：$X_s' = X_s + \sigma(\psi_w(X_s)) \odot \psi_r(R_s)$，其中$\sigma$为~Sigmoid、$\psi_w$为权重生成分支、$\psi_r$为轻量嵌入网络。这一设计使跨阶段信息传递从隐式变为显式：高置信且残差小的区域保持稳定，低置信或残差大的区域则被主动修正。

\textbf{置信度双向路由损失}——总体损失$\mathcal{L} = \sum_{s=0}^{2} w_s(\lambda_{\text{geo}}\mathcal{L}_{\text{geo}}^{(s)} + \lambda_{\text{app}}\mathcal{L}_{\text{app}}^{(s)})$通过置信度掩码实现区域级监督分配：低置信区域（$\text{Mask}_{\text{geo}}^{(s)} = \text{Mask}_{\text{gt}}^{(s)} \cdot \mathbf{1}(C_s < \tau)$）由几何损失主导，采用多尺度~Pearson~相关性同时约束局部表面起伏和全局趋势；高置信区域（$\text{Mask}_{\text{app}}^{(s)} = \mathbf{1}(C_s > \tau)$）由外观损失主导，包含~RGB~均方误差、SSIM~和~LPIPS~三项——几何修正与外观优化在三个级联阶段中互补协同。

\subsection{可泛化与多场景重建}
\label{ss:20_2}

本节讨论~2~条技术路线：（1）可泛化与高效方法；（2）城市与航天器重建。

\subsubsection{可泛化与高效方法}

本路线介绍前馈式可泛化卫星重建和大规模遥感场景的高效重建方法。

逐场景优化的卫星~3DGS~方法虽能获得高精度，但训练时间长且无法泛化到新场景。本节介绍两类改进方法：一类实现前馈式可泛化卫星重建，另一类通过频谱解耦等策略提升大规模遥感场景的重建效率。

SkySplat（2025.08）~\cite{2508_09479} 修改了可泛化~3DGS~的相机模型与训练监督环节——用卫星专用的~RPC~投影替代标准针孔相机假设，同时用相对高度监督和瞬态掩码替代纯~RGB~光度损失——高斯基元的前馈预测和体渲染保持不变。本方法比逐场景优化方法~EOGS~快~86$\times$~且精度更高，在~DFC19~数据集上将~MAE~从~13.18 m~降至~1.80 m。

\paragraph{动机与~Baseline。}标准可泛化~3DGS（如~MVSplat、HiSplat）假设针孔相机模型，即通过内参矩阵将~3D~点线性投影到~2D~像素平面。然而卫星采用推扫式（pushbroom）成像，每条扫描线的投影中心都在变化，需用~RPC~模型（一组三次有理多项式）建模像方与物方的映射关系。将~RPC~强行拟合为针孔模型会引入误差，且该误差随图像尺寸增长：在~256$\times$256~裁剪块上拟合误差为~0.27~像素，在~2048$\times$2048~时增至~2.76~像素。此外，多时相卫星影像存在三大独特挑战：（1）瞬态物体（移动车辆、云层遮挡）在不同时间采集的影像间产生不一致；（2）辐射不一致——光照和大气条件变化导致同名点亮度漂移；（3）稀疏视角——仅~3~张影像时几何约束严重不足。

\paragraph{RPC 引导的代价体构建与高度估计。}直觉上，若知道每个像素对应地面的精确高度，就能通过~RPC~反投影获得其~3D~坐标，进而直接生成高斯基元。SkySplat~的管线分四步实现这一思路。首先，用~Transformer~编码器提取多视角~2D~特征$\{F_i\}_{i=1}^N$。然后，通过可微~RPC~变换构建代价体：对参考视角的每个像素$(u_i, v_i)$，在预设高程范围内均匀采样~$M=64$~个高度候选$\{h_m\}_{m=1}^M$，利用~RPC~反投影得到~3D~坐标：

\begin{equation}
(\text{Lat}_{i}^{m},\;\text{Lon}_{i}^{m},\;\text{Hei}_{i}^{m}) = \text{RPC}_{\text{ref}}^{-1}(u_i, v_i, h_m)
\end{equation}

再将这些~3D~点通过源视角~RPC~正投影映射到源特征图上，通过插值获取源特征并变换回参考视角——直觉上，若高度假设$h_m$接近真实高程，多视角间的特征应当高度一致；否则特征不匹配。对变换后的多视角特征取方差，即得到代价体$\{C_i\}$——方差越小说明高度候选越准确。接着，对代价体沿高程维度做~soft argmin~操作，输出逐像素高度图$\{\hat{h}_i\}_{i=1}^N$。

\paragraph{高斯基元生成。}有了估计高度，通过~RPC~反投影直接计算每个像素对应的~3D~中心$\mu_{3D}$。然后用~2D U-Net~预测其余高斯参数——缩放因子$S$、旋转四元数$R$、球谐系数$C$和不透明度$\alpha$——组成完整的高斯基元集合$G = \{\mu_{3D}, S, R, C, \alpha\}$。值得注意的是，训练阶段（256$\times$256~小块）用针孔近似计算旋转四元数，推理阶段则直接使用精确的~RPC~坐标，避免近似误差累积。

\paragraph{跨自一致性模块（CSCM）。}多时相卫星影像中的瞬态物体（如移动车辆、云层）在不同视角中位置不同，若不做处理，这些区域会在训练中提供矛盾的光度监督信号，误导高斯基元学习。CSCM~的思路是：用语义特征的跨视角一致性自动检测并屏蔽瞬态区域。具体地，用~DINOv2（Distillation with NO labels v2，自监督视觉Transformer） + FeatUp~提取光度稳定的语义特征图$\{feat_i\}$，通过~RPC~变换建立跨视角特征对应关系后，计算跨视角置信度：

\begin{equation}
Q_{cv} = \max\!\left(2 \cdot \cos(\text{feat}_i,\;\text{feat}_{j \rightarrow i}) - 1,\; 0\right)
\end{equation}

其中余弦相似度低于~0.5~的位置置零——直觉上，跨视角特征相似度低意味着该像素在某个视角中被瞬态物体遮挡。对于几何过滤产生的无效区域，用自视角置信度$Q_{sv}$（参考特征与渲染特征的相似度）补偿——由于$Q_{sv}$通常低于$Q_{cv}$，需用$Q_{cv}$有效区域的均值比率校准其尺度后再替换无效区域，最终在阈值$\tau=0.2$处二值化为掩码$M$——被屏蔽区域的~RGB~损失梯度不再传播。该模块在~35k~次迭代后激活（此时高度估计已足够可靠），过早激活会因不准确的几何导致误屏蔽。

\paragraph{辐射鲁棒的自监督训练。}多时相影像间的光照和大气差异使纯~RGB~光度损失容易陷入局部最优。SkySplat~引入相对高度监督作为辅助几何信号——用~Depth Anything V2~获取单目相对高度图$H$，通过~Pearson~相关损失监督预测高度：

\begin{equation}
\mathcal{L}_{\text{hei}} = \frac{\text{Cov}(H,\;\hat{h})}{\sqrt{\text{Var}(H) \cdot \text{Var}(\hat{h})}}
\end{equation}

Pearson~相关系数衡量两个分布的形状相似性，天然免疫尺度偏移和光照变化，因此无需显式的尺度对齐步骤。RGB~损失同样受掩码保护——仅作用于稳定区域（$M=1$）：

\begin{equation}
\mathcal{L}_{\text{rgb}} = M \odot \text{LPIPS}(I_{\text{render}},\;I_{\text{gt}}) + M \odot \|I_{\text{render}} - I_{\text{gt}}\|_2^2
\end{equation}

其中$I_{\text{render}}$由标准~3DGS~的~$\alpha$-混合~渲染生成，$\odot$表示逐元素乘法——掩码$M$将瞬态区域的梯度归零，防止矛盾信号干扰高斯基元学习。最终训练损失为$\mathcal{L} = \mathcal{L}_{\text{rgb}} + \mathcal{L}_{\text{hei}}$。整个过程无需真值高度图，实现完全自监督。

\paragraph{多视角一致性聚合。}推理阶段，SkySplat~用多视角一致性策略过滤噪声点：将参考视角的预测点~RPC~投影到源视角、采样源视角高度后再~RPC~反投影回参考视角，计算像素重投影误差$\delta_p = \|p - p'\|_2$和高度相对误差$\delta_h = |\hat{h}_i(p) - \hat{h}_i(p')|/|\hat{h}_i(p)|$。仅保留$\delta_p < 3$~像素且$\delta_h < 0.2$~的可靠点，再正交投影到~2D~网格——每个网格单元保留所有分配点的最大高度——生成最终~DSM。这一数量级的加速（3.19 s vs. EOGS 的~4.60 min）得益于前馈推理架构替代了逐场景迭代优化。SkySplat~的局限性在于：作为前馈方法，其高斯基元由网络一次性生成，无法在测试时通过密度控制适应场景的局部复杂度差异。

与~EOGS~聚焦于相机模型和光照建模的系统适配不同，PSRGS（2025.03）~\cite{2503_00848} 聚焦于大规模遥感场景中~3DGS~自适应密度控制的固有缺陷：3DGS~对高频区域（建筑边缘）和低频区域（平坦地面）施加相同的密度控制阈值，而~SfM~点云在大场景中天然稀疏，导致高频区域基元不足；同时~3DGS~的平滑倾向使其生成短而宽的大型高斯椭球体，产生梯度伪影和浮游物。更关键的是，3DGS~同时优化几何与纹理时，部分纹理损失被几何损失吸收，导致高斯定位错误并持续分裂，最终恶化内存占用。

PSRGS~的核心改动是将密度控制从全局一致改为频谱引导的渐进管线——先修几何、再补纹理，分三个阶段执行：

（1）\textbf{频谱残差显著性图解耦}——对训练图像做~Fourier~变换，取对数幅度谱后用可学习高斯滤波器平滑得到残差（滤波器标准差$\sigma = \text{MLP}(S)$，其中$S$为频谱特征，与场景几何联合训练），结合原始相位谱重建显著性图，自动将场景划分为低频几何区域和高频纹理区域。滤波核为$H(x,y;S) = \frac{1}{2\pi\,\text{MLP}(S)^2}\exp\!\left(-\frac{x^2+y^2}{2\,\text{MLP}(S)^2}\right)$。

（2）\textbf{低频几何约束}——在恢复纹理细节之前先修正场景几何。对低频区域施加较高密度控制阈值（抑制多余基元），通过~$\alpha$-混合~渲染深度图$\hat{d} = \sum_{i=1}^{N}T_i\alpha_i d_i$和法向图（取协方差矩阵最短轴方向为法向），引入深度感知损失$\mathcal{L}_{\text{dep}} = \sum_i \mathcal{L}_{\text{Hu}}(\hat{d}_i, d_i^{*}) + \lambda_3\sum_{i,j}|\hat{d}_i - \hat{d}_j|^2$联合~Huber~损失与深度平滑损失约束几何正确性，其中$\lambda_3 = 0.01$。

（3）\textbf{高频纹理恢复}——同时渲染约束后的场景几何和原始未解耦场景，提取两者的梯度特征图并计算相关系数与全变分损失，确保约束后几何与原始几何对齐；然后以较低密度阈值沿梯度方向分裂和克隆高斯椭球体，保证高频区域拥有足够基元。最后将致密化后的渲染图像输入预训练网络计算纹理感知损失，进一步提升细节纹理质量。

设计动机在于：传统~3DGS~（包括~EOGS）同时优化几何和纹理，几何损失和纹理损失互相干扰——纹理误差驱动高斯基元在错误位置分裂，几何误差又反过来恶化纹理拟合。PSRGS~通过频谱解耦将几何与纹理优化在时间上分开，纹理损失不再传播到几何误差中，且不需要额外几何先验（如~SfM~稠密重建或深度估计网络）。实验中，PSRGS~在~Mip-NeRF 360~数据集上与~3DGS~使用相同超参数（15k~次后停止致密化、30k~次终止训练），在小场景中性能与其他~3DGS~变体相当，但在大规模遥感场景中在高频细节恢复和几何精度上明显优于对比方法。

\subsubsection{城市与航天器重建}

本路线介绍卫星到地面的极端视角外推、航天器~3D~重建和城市大规模重建等方法。

城市和航天器的~3D~重建是卫星遥感中两个典型且互补的应用场景：城市重建需要从天底俯视影像恢复建筑立面等被遮挡结构，航天器重建则需要从地面或轨道影像重建无纹理的金属表面。本节介绍~3~种代表性方法，分别从卫星到地面的极端视角外推、无纹理航天器重建以及大规模城市场景的~TDOM~生成等方向展开讨论。

\paragraph{卫星到地面的极端视角外推}
本方法修改了~3DGS~的\textbf{场景表示与纹理建模}环节——用~Z-Monotonic SDF（2.5D 高度图）替代自由形式的~3D~高斯基元表示几何，用扩散模型驱动的迭代纹理恢复替代直接反投影——基于可微分光栅化的渲染管线保持不变。

\paragraph{问题：为何卫星影像合成地面视角极其困难？}
从卫星影像到地面视角的视角差距接近~90\textdegree{}（天底俯视到地面平视），属于极端新视角外推。在这一条件下，标准~3DGS~面临两个根本性困境。第一，几何表示困境：标准~MVS~算法只能重建屋顶和地面的密集点云，建筑立面几乎无点云——因为卫星看不到立面。此时~3DGS~的优化器面对高度病态的条件，自由形式的高斯基元会产生"收缩包裹"（shrink-wrapped）几何，完全无法捕捉真实的城市结构。第二，纹理困境：即使获得了正确几何，卫星影像的分辨率不足以提供立面纹理细节，直接反投影卫星图像会将模糊、锯齿和投影伪影"烘焙"进纹理，合成的地面视角无法令人信服。

\paragraph{核心改动一：Z-Monotonic SDF 几何表示。}
为了解决几何表示困境，本方法引入了一个强城市结构先验：将城市建模为~2.5D~高度图（即~3D~表面可由地面平面上的单值高度函数~$z = f(x, y)$~定义）。直觉上，这意味着城市中的每个~$(x, y)$~位置只有一个高度值——建筑从地面"长出来"，不会有桥梁等非单调结构。这一表示与卫星的自上而下观测特性和城市建筑以竖直结构为主的特点天然匹配。

具体实现上，一种朴素做法是先从卫星影像用~MVS（Multi-View Stereo，多视角立体）~提取点云，再直接将其转化为~2.5D~高度图（即遥感文献中的~DSM，Digital Surface Model），然后填充~3D~体素网格并用~Marching Cubes~提取表面——然而由于~MVS~点云稀疏且含噪声，这种直接转化会在建筑边缘产生严重的混叠（aliasing）和"阶梯状"伪影。本方法不采用这种朴素管线，而是优化一个带有~Z~轴单调约束的~SDF（有符号距离函数）场~$s(\mathbf{x})$，强制其值沿竖直方向非递减：

\begin{equation}
    \frac{\partial s(x, y, z)}{\partial z} \ge 0 \quad \forall \mathbf{x} = (x, y, z).
\end{equation}

直觉上，这个约束意味着~SDF~值"从下往上只能增大或不变"——在连续表面（如屋顶）处，方程~$s = 0$~有唯一解，定义了高度值；在建筑边缘处，函数退化为竖直平台，即~$s(x, y, z) = 0$~对所有~$z \in [z_{\text{ground}}, z_{\text{roof}}]$~成立，程序化地定义了完美的竖直立面。通过这一约束，原本高度病态的立面重建问题被转化为一个约束良好的优化问题。

在实现层面，Z-Monotonic SDF~通过在~$256 \times 256$~的~2D~网格上参数化单调曲线来实现。直觉上，每个网格单元存储一个"台阶高度"参数，整条单调曲线就像一组~tanh~阶梯函数的叠加——每个~tanh~负责一个高度突变，组合起来就能逼近任意的单调~SDF~轮廓。对于查询点~$\mathbf{p} = (x, y, z)$，SDF~值通过对邻域内多条~tanh~基函数曲线的输出进行加权插值得到：

\begin{equation}
    f(z; x, y) = \sum_{j=1}^{n} w_j \cdot \tanh(k \cdot (z - h_j)),
\end{equation}

其中~$h_j$~为可学习的竖直偏移参数，$k$~为控制曲线陡峭度的超参数，$w_j$~为基于空间邻近度的插值权重。几何优化通过最小化~MVS~点云到提取网格的~Z~轴距离（而非~3D~距离，因为水平定位已由卫星~RPC~模型提供），配合~Laplacian~正则化和法向一致性损失：

\begin{equation}
    \mathcal{L}_{\text{geo}} = \sum_{\mathbf{p} \in P} \min_{\mathbf{m} \in M} \|\mathbf{p}_z - \mathbf{m}^*(\mathbf{p})_z\|_1 + \lambda_{\text{Lap}} \mathcal{L}_{\text{Lap}} + \lambda_{\text{Nrm}} \mathcal{L}_{\text{Nrm}},
\end{equation}

其中~$\mathbf{m}^*(\mathbf{p})$~是在~XY~平面上与~$\mathbf{p}$~最近的网格点。使用~Z~轴距离而非~3D~距离是刻意的设计选择：卫星影像的~RPC~相机模型已提供高精度的水平定位，需要优化的主要是高度。几何优化使用~Adam~优化器，学习率~$0.01$，通过可微分等值面提取技术（2023）~\cite{shen2023flexible}从~SDF~场中提取水密网格。

\paragraph{核心改动二：生成式纹理恢复。}
在几何建立后，本方法将外观建模与几何完全解耦——纹理不再由高斯基元的颜色属性表示，而是作为~UV~纹理贴图绑定到网格上。初始纹理通过可微分渲染器反投影卫星图像生成，但直接反投影的结果模糊且充满伪影。为此，本方法引入一个图像恢复网络~$D$（基于~FLUX-Schnell~扩散模型微调，在~100{,}000~对航拍图像上训练），将退化渲染结果映射为清晰纹理。恢复网络通过感知损失和~Charbonnier~保真损失训练：

\begin{equation}
    \mathcal{L}_{\text{restorer}} = \mathcal{L}_{\text{LPIPS}}(\hat{I}, I_{\text{high}}) + \lambda_{\text{CHAR}} \mathcal{L}_{\text{CHAR}}(\hat{I}, I_{\text{high}}).
\end{equation}

这里使用~LPIPS~而非纯~MSE~损失，是因为卫星影像的模糊和降质不仅需要"去噪"，更需要"生成"卫星从未观测到的立面细节——感知损失能引导网络产生视觉上合理的细节，而不仅是像素级忠实。恢复网络训练~10{,}000~次迭代（batch size~$96$，$\lambda_{\text{CHAR}} = 1$），在配对的城市~3D~场景数据集上学习从退化图像到高质量目标的映射。

纹理优化采用迭代精炼策略：每轮迭代先从当前纹理贴图渲染模拟~UAV（Unmanned Aerial Vehicle，无人机）~视角的近景图像（退化渲染），再由恢复网络增强为清晰伪真值，然后用这些增强图像作为监督信号更新纹理贴图。模拟视角的采样方式为：在网格包围盒扩展~100\,m~的区域内以~150\,m~间距均匀采样，模拟飞行高度~450\,m、俯仰角~$45^{\circ}$、四个主方向的~UAV~观测。这一"渲染$\to$增强$\to$更新"的反馈循环使纹理逐步逼近照片级真实——每一轮的精炼纹理为下一轮提供更高质量的输入，实现"自举式"渐进提升。具体参数上，初始纹理优化~100~个~epoch（$\lambda_{\text{MSE}} = 0.8$，$\lambda_{\text{SSIM}} = 0.2$），迭代精炼共~2~轮，每轮生成~$2048 \times 2048$~分辨率的目标图像，优化~20~个~epoch。整个流程在单块~NVIDIA A6000 GPU~上仅需约~1.5~小时处理~1 km$^2$~城市区域（几何优化~0.5~小时+外观精炼~1~小时），最终在~4 km$^2$~真实城市区域上验证了从卫星视角合成可信地面视角的能力。

\paragraph{设计动机总结。}
本方法的核心设计决策——用~2.5D~高度图替代~3DGS~的自由高斯基元表示——本质上是在表达能力和鲁棒性之间做取舍。2.5D~表示牺牲了对桥梁、悬挑等非单调结构的建模能力，但换来了对几何歧义性的强力正则化：当卫星看不到立面时，自由形式表示会崩溃，而~2.5D~表示通过单调性约束"补全"了立面。类似地，用生成式恢复网络替代直接反投影，本质上是利用扩散模型从大规模航拍数据中学习的城市外观先验，"想象"卫星从未看到的细节。

Park and D'Amico（2025.12）~\cite{2512_23998} 从同一卫星/航天器重建领域出发，但将关注点从"地面场景的几何与纹理"转向"太空环境中~3DGS~对照明变化的适应性"。上一方法（From Orbit to Ground）假设卫星图像之间存在相对一致的成像条件，因此将精力集中于几何表示与纹理生成；而~RPO（近距离交会操作）场景面临截然不同的光照挑战——在低地球轨道上，太阳方向随轨道位置和航天器姿态快速变化，方向性阳光在目标航天器（RSO）表面投射极端且剧烈变化的阴影。标准的~in-the-wild~方法（如~WildGaussians）为每张图像分配可学习嵌入向量~$\boldsymbol{e}_j$，通过外观~MLP~将嵌入映射为视角相关颜色~$\hat{c}_i = A(f_i, \boldsymbol{e}_j)$——但这要求对新图像执行测试时优化以求解~$\boldsymbol{e}_j$，在星载计算资源受限的~RPO~场景中不可行。

\paragraph{核心改动：用太阳方向向量替代逐图像外观嵌入。}
本方法的关键洞察是：在太空中，太阳是唯一的主导光源，而每颗自主导航卫星的姿态轨道控制系统（AOCS）始终维护着高精度的太阳方向估计。因此，无需为每张图像学习一个"黑箱"嵌入，可以直接用已知的太阳方向向量~$\boldsymbol{S}_j^{\text{sun}}$~作为外观条件输入。具体实现上，每个高斯配备~72~维可训练特征向量~$\boldsymbol{f}_i \in \mathbb{R}^{72}$（而非原始~3DGS~的球谐系数），经位置编码后的太阳方向向量和相机视角向量一起输入外观~MLP$A$（3~层~256~维隐藏层）预测颜色：

\begin{equation}
    c_i = A\bigl(\boldsymbol{f}_i,\; PE(\boldsymbol{S}_j^{\text{sun}}),\; PE(\boldsymbol{d}_{ij}^{\text{view}})\bigr),
\end{equation}

其中位置编码~$PE_L(\boldsymbol{p})$~采用~NeRF~式多频率正弦/余弦映射。这一替换的本质在于：将"需要优化的自由参数"替换为"物理已知的观测量"，从而消除了测试时优化需求，使~3DGS~能在星载平台上实现即插即用的外观预测。

\paragraph{扩展一：阴影~splatting。}
仅用太阳方向预测颜色仍不足以捕捉太空中的自遮挡阴影现象——方向性阳光在~RSO~表面投射出极端锐利的阴影，这是~MLP~难以从数据中隐式学习的。本方法采用~GS$^3$~的阴影~splatting~机制：将高斯沿光源方向~splat（等效于将相机置于太阳位置），通过累积不透明度计算每个高斯的可见度~$V_i \in [0, 1]$（1~表示被照亮，0~表示被遮挡）。该可见度经一个小型精炼~MLP$R$（3~层~32~维隐藏层）进一步优化后生成阴影图像，与主渲染图像逐像素相乘：

\begin{equation}
    V_i^{*} = R\bigl(V_i,\; PE(\boldsymbol{S}_j^{\text{sun}}),\; \gamma(\boldsymbol{\mu}_i),\; \boldsymbol{s}_i\bigr),
\end{equation}

其中~$\gamma(\boldsymbol{\mu}_i)$~为高斯中心位置编码，$\boldsymbol{s}_i \in \mathbb{R}^6$~为可训练阴影潜变量。这一设计的精妙之处在于复用了~3DGS~已有的光栅化管线——无需额外的光线行进或体素化，仅通过改变"相机"位置即可获得遮挡信息。

\paragraph{扩展二：等向尺度正则化。}
本方法引入等向尺度正则化损失，鼓励高斯基元趋于球形：

\begin{equation}
    \mathcal{L}_s = \sum_{i=1}^{|G|} \|\boldsymbol{S}_i - \bar{S}_i \mathbf{1}\|_1,
\end{equation}

其中~$\boldsymbol{S}_i$~为高斯的尺度参数向量，$\bar{S}_i$~为其均值。这一正则化最初受~GS-SLAM~启发，旨在为~SLAM~中的持续跟踪提供稳定特征——但其与阴影~splatting~的协同效应更为关键：规则球形高斯使每个基元的遮挡范围更精确、更局部化，从而渲染出更锐利的阴影轮廓。实验表明，加入~$\mathcal{L}_s$~后高斯变得更加球化和细粒度，但也带来高斯数量增加的副作用。总训练损失为~$\mathcal{L} = \mathcal{L}_{\text{3DGS}} + \lambda_s \mathcal{L}_s$，其中~$\mathcal{L}_{\text{3DGS}} = (1 - \lambda_{\text{D-SSIM}}) \mathcal{L}_1 + \lambda_{\text{D-SSIM}} \mathcal{L}_{\text{D-SSIM}}$，$\lambda_{\text{D-SSIM}} = 0.2$。在关键帧管理上，采用~DSO~式启发策略：当新帧与最近关键帧的视角偏差超过~$\theta_{\min} = 10^{\circ}$~时添加为新关键帧，窗口最多保留~$W = 10$~帧（增加~$W$~可改善性能但增加内存开销）。整个流程在~Cassini、Rosetta、Starlink~三种卫星模型上验证——完整模型（外观~MLP + 阴影~splatting + 等向正则化）在所有指标上超越纯~3DGS~基线，其中阴影~splatting~带来最显著的单项提升。

与上述方法利用~3DGS~进行太空场景重建和外观建模的思路不同，SatSplatYOLO（2024.06）~\cite{2406_02533} 将~3DGS~的渲染能力重新定位为一个"虚拟观测生成器"——其核心目标不是重建本身，而是通过渲染合成视角为下游航天器部件检测任务提供数据增强。

\paragraph{问题：单视角检测的固有瓶颈。}
在近距离交会操作（RPO）中，基于~CNN（Convolutional Neural Network，卷积神经网络）~的航天器部件检测器（如~YOLOv5）面临一个根本矛盾：部署在星载边缘硬件上只能使用轻量级模型，而轻量模型在未知卫星上的泛化能力不足——尤其是太阳能板、天线、推进器、本体四类部件的检测极易产生假阳性（将本体误检为天线）和假阴性（遗漏小部件）。先前研究已证明，使用多个物理观测器能显著改善检测结果，但在实际~RPO~中让多个航天器近距离围绕目标飞行既不经济也不安全。

\paragraph{核心改动：3DGS 驱动的虚拟视角集成检测。}
本方法保留标准~3DGS~的~SfM~初始化和可微分光栅化管线不变（不做任何场景表示层面的修改），而是在此基础上设计了两个新增模块。

第一，系统性虚拟相机生成。给定~$n$~个真实相机位姿~$T_i$，首先通过最小二乘求解所有相机视线方向到公共注意力中心点~$\mathbf{p}$~的最优估计：

\begin{equation}
    \mathbf{p}^* = \arg\min_{\mathbf{x}} \|A\mathbf{x} - C\|_2^2,
\end{equation}

其中~$C$~为所有相机位置拼接的列向量，$A$~为包含视线方向向量的分块矩阵。然后围绕每个真实相机位置生成~$m$~个虚拟位姿——支持圆形采样（在垂直于视线的平面上以半径~$r$~等间隔采样~$c^*_{ij} = c_i + r(\cos\theta_j \mathbf{u}_i + \sin\theta_j \mathbf{w}_i)$）和球形采样（使用~Fibonacci~格点半均匀覆盖球面），虚拟相机始终指向注意力中心~$\mathbf{p}$。每个真实视角生成~$64$~个虚拟渲染，共~$84 \times 64 = 5376$~张合成图像。

第二，基于检测一致性的集成策略。将真实图像和~$64$~张合成渲染分别送入~YOLOv5-s~（选用~small~架构以适配星载硬件的实时推理需求）进行部件检测。集成过程分为四步：（1）将合成视角中~IoU~高于阈值的检测框聚类为组；（2）计算每组的多数类预测和纯度（purity，即与多数类一致的检测比例）；（3）将纯度低于阈值的组过滤，将同类别且~IoU~重叠的组合并；（4）执行平移校正——由于~SfM~和注意力中心近似引入的系统性偏差，合成视角检测框相对真实视角存在可预测的平移偏移，通过将检测组均值框映射到真实图像上高置信度~YOLOv5~预测的位置来修正。最终，若合成检测组与真实检测在空间上重合，则用~$(\text{confidence}, \text{purity})$~加权平均更新检测框和类别；若不重合，则补充为新检测（修正假阴性）；同时移除无合成检测组对应的真实预测（抑制假阳性）。

\paragraph{设计动机。}
本方法的核心洞察在于：3DGS~的实时渲染能力（相比~NeRF~的逐射线采样）使得大规模虚拟视角生成在实际可行的计算预算内成为可能，从而将"多观测器协作"的思想从物理空间转移到虚拟空间。值得注意的是，SatSplatYOLO~在已知卫星上的~mAP@0.5~为~0.825（与基线~YOLOv5~持平，因为已知卫星的检测已足够准确），但在未知卫星上~mAP@0.5~从~0.515~提升至~0.610（提升~18.4\%），验证了集成策略对泛化场景的增益。3DGS~模型本身的渲染质量达到~SSIM~$= 0.9213$、PSNR~$= 25.52$~dB、LPIPS~$= 0.0796$，为下游检测提供了高质量的虚拟观测基础。

Park and D'Amico~的~RPO~场景虽然光照条件苛刻，但观测距离近（数米到数十米）、姿态已知、成像质量高。Chang et al.（2024.04）~\cite{2404_18394}~转向了观测条件的另一个极端——从业余望远镜的地面观测图像重建~3D~卫星模型（包括天宫空间站和国际空间站），观测距离数百公里以上，且相机位姿完全未知。与锚点方法依赖~AOCS~提供的已知太阳方向和姿态不同，本方法必须从零开始估计相机位姿，同时应对大气湍流导致的严重模糊和畸变。

\paragraph{核心差异一：从~RPO~近景成像到地面远距成像，需要图像预处理。}
锚点方法在~RPO~场景中成像距离近，无需图像恢复。而地面望远镜观测面临大气湍流的系统性干扰——单帧图像严重模糊且畸变。本方法为此设计了三级混合预处理流水线：第一级，Lucky Imaging——从高速短曝光序列中按清晰度和对比度指标统计选取最锐利的~12\%~帧，经质心对齐和逐块归一化互相关校正子像素抖动后叠加，将零均值随机噪声降低~$1/|S|$~倍（$|S|$~为选中帧数）；第二级，6~级~Gaussian~小波分解多尺度锐化——细尺度（$256 \times 256$）到粗尺度（$8 \times 8$）的增益依次为~$(1.1, 1.1, 1.1, 1.1, 1.1, 5.0)$，最粗尺度的增益大幅提高以恢复太阳翼等低对比度扩展结构；第三级，Video Restoration Transformer（VRT）去噪——利用自注意力机制和时序一致性进一步抑制锐化引入的伪影，保留结构细节。

\paragraph{核心差异二：从已知姿态到联合位姿估计与~3D~重建。}
锚点方法中相机姿态由航天器~AOCS~系统直接提供，不涉及位姿估计。而地面望远镜的观测位姿完全未知，本方法必须将位姿估计与~3D~重建耦合为联合优化问题。初始化阶段采用改进的~SfM：用~SuperPoint+SuperGlue~在望远镜图像的稀疏纹理和畸变条件下提取并匹配特征，采用正交投影模型（而非透视模型）初始化前三帧——因为数百公里的观测距离使透视效应极弱，正交投影更鲁棒——然后切换回透视模型逐帧加入新图像，每加入一帧执行全局光束法平差防止误差累积，每五帧进行三角化补充新特征点。

初始化完成后，核心重建阶段采用受控~3DGS~与~Branch-and-Bound（BnB）位姿精炼的交替迭代。与锚点方法中高斯自由优化不同，本方法采用受控增长策略：早期训练阶段固定高斯数量，仅优化位置、协方差、不透明度和球谐系数（最高三阶），强制每个高斯覆盖更大区域以捕捉粗糙几何，防止在位姿不精确时过拟合高频伪影；经过设定轮数后自动克隆分裂增加高斯数量，逐步捕获精细细节。与此同时，每~1000~轮通过~BnB~搜索精炼相机位姿——生成~200~个扰动候选位姿，旋转扰动通过~Rodrigues~公式$R_{\text{noise}} = I + (\sin\theta)\,[\mathbf{n}]_\times + (1 - \cos\theta)\,[\mathbf{n}]_\times^2$采样，平移扰动从零均值高斯分布$T_{\text{noise}} \sim \mathcal{N}(0, \sigma_T^2)$采样——用光度一致性损失评估并选择最优位姿，随后将搜索范围缩小（平移至~$1/4$、旋转至~$1/2$），确保~10~轮内稳定收敛。此外，由于地面望远镜图像含背景噪声，重建的点云中可能包含杂散点，本方法增加了一步统计滤波后处理：计算每个高斯到其~$k$~近邻的平均距离，移除超过全局均值加一个标准差的点，然后以滤波后的高斯集合恢复训练~500~轮以修复不连续并提升重建质量。

\paragraph{设计动机。}
锚点方法~Park and D'Amico~的核心设计思路是"在已知的良好条件下增强~3DGS~的表达能力"——已知姿态下用太阳方向驱动外观、用阴影~splatting~增强几何感知。本方法~Chang et al.~的设计动机则恰恰相反："在极端恶劣的观测条件下先稳定优化、再逐步精细化"——图像预处理消除系统性退化，受控高斯增长防止早期过拟合，BnB~搜索在离散空间中可靠收敛位姿，统计滤波清除噪声杂点。整个流水线无需逐场景调参，在合成数据、CSS（火箭上面级）和~ISS（国际空间站）上使用完全相同的超参数集合，重建质量超越了基于~NeRF~的方法。

\begin{takeaway}
卫星 3D 重建面临极端观测条件（弱纹理、未知位姿、极端光照）的挑战：太阳方向驱动的外观建模和阴影感知 splatting 在已知条件下增强表达能力，受控高斯增长和 BnB 位姿搜索则在未知条件下保证稳定收敛。
\end{takeaway}

%=====================================================================

\subsection{近期方法脉络：卫星与RPC高斯渲染、Satellite photogrammetry with 2D Gaussian Splatting、Generative refinement for satellite Gaussian Splatting}

这一组近期工作延续 Gaussian Splatting、2DGS、Satellite Gaussian Splatting，主要把改动集中在扩展、适配等环节。下面只保留相对既有路线的变化，便于和前文方法直接对照。

\paragraph{RPC-GS}
RPC-GS（2026.06）~\cite{2606_06690} 延续 Gaussian Splatting 在“卫星与RPC高斯渲染”上的基本设定，主要改变是：RPC-GS：将Gaussian均值和协方差直接通过Rational Polynomial Camera模型投影，并推导RPC协方差映射。3DGS 的显式高斯表示和可微渲染能力保持为基础。
其设计动机是：把GS投影从透视相机扩展到卫星RPC几何，减少pushbroom传感器近似带来的重建误差。
具体而言，方法构建从适合splatting的场景坐标到地理坐标再到RPC图像坐标的变换链，并数值稳定地映射高斯协方差矩阵。
适用边界是：RPC投影和协方差传播计算复杂；依赖卫星元数据质量。

\paragraph{SatSplat}
SatSplat（2026.06）~\cite{2606_28581} 延续 2DGS 在“Satellite photogrammetry with 2D Gaussian Splatting”上的基本设定，主要改变是：将 2DGS 适配到卫星摄影测量，并进行在线相机调整。3DGS 的显式高斯表示和可微渲染能力保持为基础。
其设计动机是：针对卫星多日期、高空小交会角和光照差异导致的重建退化，提出适合卫星相机模型的 Gaussian Splatting 框架。
具体而言，以 2DGS 为基础，使用 affine camera approximation，并学习最小 delta 参数化进行 in-splat camera refinement，从而缓解卫星相机和几何误差。
适用边界是：卫星观测顶视角限制和多时相外观差异仍会影响立面和遮挡区域；相机近似适用范围需注意。

\paragraph{SatSplatDiff}
SatSplatDiff（2026.06）~\cite{2606_27223} 延续 Satellite Gaussian Splatting 在“Generative refinement for satellite Gaussian Splatting”上的基本设定，主要改变是：在卫星 GS 中引入几何保持的生成式渲染监督细化。3DGS 的显式高斯表示和可微渲染能力保持为基础。
其设计动机是：利用生成先验改善卫星 GS 的视觉保真度，同时约束独立视图细化可能造成的幻觉和几何不一致。
具体而言，在 Gaussian 渲染训练中迭代使用生成模型细化监督图像，并通过几何保持机制减少跨视图 hallucination 对建筑几何的破坏。
适用边界是：生成式细化可能引入不真实细节；卫星多日期光照、稀疏顶视角和立面不可见仍是挑战。
\section{航空影像重建}
\label{sec:20_2}

无人机（UAV，Unmanned Aerial Vehicle）航拍影像为~3DGS~提供了中等距离、高分辨率的观测条件，但仍面临多时相光照不一致、大规模场景分块处理和语义理解等挑战。

\subsection{航拍重建方法}
\label{ss:20_3}

本节讨论~4~条技术路线：（1）实时重建与~LiDAR~融合；（2）语义重建与物理仿真；（3）高精度表面重建；（4）数据集与评测基准。

\subsubsection{实时重建与LiDAR融合}

本路线介绍无人机端到端实时重建管线和~LiDAR~引导的高精度航拍重建方法。

\paragraph{Maikos et al.：端到端实时重建管线}

Maikos et al.（2026.02）~\cite{2602_20342} 提出端到端管线将无人机视频流转换为高保真~3D~重建，集成~RTMP~实时视频流、传感器融合和~3DGS~优化，渲染质量在~4--7\%~以内接近离线高保真参考。

\textbf{改动概要。}本方法修改了~3DGS~的数据采集与优化环节——用~RTMP~实时视频流和在线增量优化替代离线图像集合的批量训练——3DGS~的高斯表示和可微光栅化渲染管线本身保持不变。读者已从第~1~章了解标准~3DGS~的训练流程：离线~SfM~生成稀疏点云$\to$初始化高斯$\to$可微光栅化优化。该方法的核心问题是：能否将这个离线闭环打通为实时管线，让无人机边飞边重建，同时保持接近离线方法的质量？

\textbf{Baseline~做法与挑战。}标准~3DGS~假设所有训练图像已预先采集完毕，通过~COLMAP~一次性完成~SfM~重建后，再在整个数据集上训练~30\,000~次迭代。这对无人机实时场景有两个根本限制：第一，无法在飞行过程中持续获得更新——无人机看到的新区域无法即时融入已有模型；第二，NeRF~类方法（Instant-NGP、Mip-NeRF~360）虽可提供新视角合成，但光线行进（ray marching）的渲染速度远低于~3DGS~的基于瓦片的光栅化，无法满足~AR（Augmented Reality，增强现实）~交互所需的帧率。

\textbf{管线设计。}该系统将无人机采集到~3D~可视化的全链路分为五个串联模块，每个模块针对实时性做了专门设计：

\begin{enumerate}
\item \textbf{RTMP~实时视频流采集}——每架无人机通过硬件加速~H.264/H.265~编码将视频压缩后经由~RTMP~（Real-Time Messaging Protocol，一种基于~TCP~的实时消息传输协议）传输至地面服务器。控制指令、遥测数据和视频数据走独立通道，避免大~I~帧阻塞飞行控制包。服务器端解码为原始帧后通过共享内存传递给后续处理管线，端到端延迟控制在数百毫秒级别。
\item \textbf{帧提取与多模态同步}——多路传感器数据（视频、IMU、GPS）统一到~IEEE 1588 PTP~时间基准。同步模块为每种模态维护独立队列，将视频帧与时间窗内最近的传感器采样对齐；缺失数据通过插值（线性度量）或积分（惯性测量）补全。为控制计算负载，帧率可均匀或自适应降低——根据运动幅度或分割动态性选择性保留关键帧，而高速内容保留完整帧序列以捕捉精细时序细节。
\item \textbf{相机位姿估计}——对齐后的传感器数据送入~SfM~管线，输出每个相机帧在全局坐标系中的~6-DoF~位姿，即~SE(3)（Special Euclidean Group，特殊欧几里得群）~变换
\begin{equation}
T_i = [R_i \mid t_i], \quad R_i \in \mathbb{R}^{3\times3},\; t_i \in \mathbb{R}^3,
\end{equation}

其中~$R_i$~为旋转矩阵，$t_i$~为平移向量。这个公式的直觉含义是：每个相机帧的空间位置由一个刚体变换（旋转+平移）唯一确定。IMU/GPS~辅助信息和语义分割在无人机抖动、弱纹理等困难场景下稳定位姿估计。
\item \textbf{3DGS~在线优化}——高斯初始化于~SfM/MVS~点云位置，通过可微瓦片光栅化器计算光度损失并以~mini-batch~方式优化参数，以最大化~GPU~利用率。训练过程采用混合精度（mixed-precision）计算加速收敛，同时引入空间正则化（spatial regularization）防止相邻高斯之间产生接缝伪影（stitching artifacts）。密度自适应控制机制（densification~在重建不充分的区域注入新高斯，pruning~移除低贡献高斯）确保模型紧凑和几何完整。

\textbf{关键区别}在于：部署阶段不从头重训练，而是仅对受新数据影响的区域执行在线优化——即仅重新优化受新帧影响的空间区域内的~splat，而非全场景重训练。模型以紧凑二进制格式（连续参数数组+空间瓦片数据）存储，支持增量写入。
\item \textbf{WebSocket~实时分发}——训练好的模型通过~WebSocket~推送至客户端。客户端可选择替换、合并或按区域选择性加载新模型，支持~Unity~引擎中的连续资产流式传输和多用户协作。
\end{enumerate}

\textbf{设计动机。}传统航空~MVS~方法（COLMAP、OpenMVS）几何精度高但无法提供实时渲染和交互式新视角合成；NeRF~方法可合成新视角但光线行进的渲染速度（通常~$<$1~FPS）远低于~AR~所需帧率。该系统选择~3DGS~作为重建核心，正是利用其两个独特优势：基于瓦片光栅化的渲染速度（实验中达到~134--154~FPS）和高斯的局部可更新性——可以仅对局部区域重新排列~splat~而不会产生~NeRF~全局更新的副作用。此外，各向异性高斯核能有效减少混叠和几何间隙，球谐函数系数可捕获视角相关光照，弥补了标准~MVS~管线缺乏外观建模的不足。

\textbf{实验结果。}在~Mip-NeRF~360、Tanks and Temples~和~Deep Blending~三个基准上，完整训练~30\,000~次迭代的配置（Ours30K）达到~SSIM~0.815--0.903、PSNR~27.21--29.41~dB、LPIPS~0.183--0.243、渲染速度~134--154~FPS，重建质量保持在离线高保真参考的~93--96\%~范围内（差距~4--7\%）。缩短至~7\,000~次迭代的配置（Ours7K）在训练时间大幅缩减的同时维持了可比的图像质量指标和更高的渲染速度。与~Instant-NGP~和~Mip-NeRF~360~的对比表明，3DGS~方法在训练和渲染速度上均显著优于~NeRF~方法，同时在视觉保真度上保持竞争力。实验硬件为~AMD Ryzen 9 9900X3D + NVIDIA RTX 4070 12~GB + 64~GB~RAM，无人机为~DJI Mini 4K。

\paragraph{ARSGaussian：LiDAR引导航拍重建}

ARSGaussian（2024.12）~\cite{2412_18380} 沿着锚点方法「无人机~$\to$~3DGS~重建」的基本范式继续前进，但指出该范式中一个被忽视的薄弱环节：锚点方法（以及大多数纯视觉~3DGS~方法）完全依赖~SfM~生成的稀疏点云作为初始化和几何基准，而航拍场景中远距离、稀疏视角的采集条件使得~SfM~的几何估计极易出错——具体表现为两类典型退化：浮游体（floater，空中飘浮的虚假高斯）和过度生长（overgrowth，高斯向观测稀疏区域无序扩张，在渲染图像中表现为蓝色圈内的过度拉伸区域）。ARSGaussian~的核心回答是：能否用~LiDAR（Light Detection and Ranging，激光雷达）~点云的精确几何作为「锚」，把~3DGS~的高斯基元约束在真实表面附近？

\textbf{改动一：LiDAR~引导的自适应致密化策略。}锚点方法在~3DGS~在线优化模块中使用了标准的~densification/pruning~机制——根据屏幕空间梯度判断哪些区域需要更多高斯、哪些高斯贡献过低。ARSGaussian~不改变这一优化管线本身，而是在致密化判断逻辑中引入~LiDAR~几何基准，将原来的单一「梯度阈值」规则扩展为三维约束：（1）\textbf{冗余高斯移除}——每~300~次迭代，用~Kd-tree~搜索每个高斯最近的~LiDAR~点，分别计算水平面距离（$d_{\text{plane}}$，由经纬度差值乘以地球半径得到）和高程距离（$d_{\text{elev}}$），若两者同时超过各自的阈值则标记为待移除；但关键细节是，如果该高斯正上方存在~LiDAR~点（即~$d_{\text{plane}} < 2\epsilon$~且高程差为正），即使超出距离阈值仍予保留——这是对航拍场景特殊性的针对性设计：机载~LiDAR~向下扫描时水平覆盖完整但垂直面常有缺失（如建筑立面），简单设单一阈值会误删合法高斯。（2）\textbf{分裂方向引导}——对于欠重建的高斯（平均梯度超过阈值$\tau_{\text{pos}}$），将高斯椭球的长轴（协方差矩阵最大特征值对应的特征向量）投影到最近~LiDAR~点的局部切平面上，以投影方向作为新生成高斯的生长方向。这样，高斯不是沿随机方向分裂，而是沿着~LiDAR~揭示的真实几何表面延展，从根本上抑制了浮游体的产生。

\textbf{改动二：LiDAR-图像像素级对齐模块。}上述致密化策略的前提是~LiDAR~点云与光学图像之间实现像素级对齐——直接投影误差可达数十像素，在高度变化剧烈的区域尤为严重。锚点方法采用标准~COLMAP SfM~估计相机位姿，使用简单针孔相机模型。ARSGaussian~则用~Brown--Conrady~畸变模型（引入二阶和四阶径向畸变系数$k_1$、$k_2$）替代针孔模型来消除遥感光学系统的复杂畸变，将~DGPS/IMU（Differential GPS / Inertial Measurement Unit）~提供的初始位姿作为先验，再通过~Colmap-PCD~框架联合优化：将~LiDAR~点云体素化后投影到图像平面，与图像特征点匹配并最小化重投影误差，经三次渐进式光束法平差得到精确相机位姿，实现~LiDAR~点云与多视角图像在~WGS84~世界坐标系中的像素级对齐。对齐后的融合点云（LiDAR~点+图像特征点）替代~SfM~稀疏点云作为~3DGS~的初始化输入。

\textbf{改动三：三重几何一致性约束。}锚点方法的光度损失（$L_1$+D-SSIM）仅约束渲染图像与真实图像的像素级一致性，不直接约束几何精度。ARSGaussian~在光度损失之上增加几何损失$L_{\text{geo}} = \alpha L_{\text{depth}} + \beta L_{\text{normal}} + \gamma L_{\text{scale}}$，三项损失均以~LiDAR~生成的稠密深度图和法向图为监督信号。由于机载~LiDAR~分辨率低于光学图像（4--8 pts/m$^2$），需要先通过~Delaunay~三角剖分插值+ACMH~多视角立体匹配迭代优化，将稀疏~LiDAR~深度传播为稠密深度图，再从深度图推导法向图。深度损失$L_{\text{depth}}$和法向损失$L_{\text{normal}}$均采用~$L_1$~距离与余弦相似度的组合形式（$\|\hat{D}(p) - \bar{D}(p)\|_1 + \|1 - \hat{D}(p)^\top \bar{D}(p)\|_1$），同时约束绝对距离和方向一致性；尺度损失$L_{\text{scale}}$约束高斯最短轴趋近于零以抑制过度膨胀。

\textbf{设计动机。}锚点方法回答了「无人机能否实时重建」，ARSGaussian~则追问「重建的几何精度够不够」——纯视觉~3DGS~在航拍场景中因远距离、稀疏视角导致几何退化是固有缺陷，而机载~LiDAR~天然提供厘米级精度的三维坐标，将两者融合可以在保持~3DGS~渲染速度优势的同时补齐几何精度短板。论文还发布了~AIR-LONGYAN~数据集（福建龙岩约~32.5~万~m$^2$，机载~LiDAR~密度~4--8 pts/m$^2$+五镜头倾斜相机~602~张图像），填补了航拍场景中缺少同步~LiDAR-光学混合数据集的空白。在~UrbanScene3D~和~AIR-LONGYAN~上，ARSGaussian~在~PSNR/SSIM/LPIPS~上均优于~VastGaussian、GaussianPro、Pixel-GS~等方法，且~LiDAR RMSE~（全场景高斯到~LiDAR~点的均方根误差）显著低于所有对比方法。

\subsubsection{语义重建与物理仿真}

本路线介绍基于~3DGS~的语义航拍理解和滑坡物理仿真方法。

\paragraph{Chick Zaouali et al.：语义航拍理解}

语义航拍重建方面，Chick Zaouali et al.（2025.05）~\cite{2505_17402} 将问题从上述~Maikos et al.~的「实时重建」转向了「重建后语义理解」——Maikos et al.~的管线输出的是纯几何与外观模型，缺少对场景内容的语义解读能力；而实际航拍巡检（基础设施检测、结构评估）恰恰需要知道「重建出的场景里有什么、在哪里」。

\textbf{方法设计。}本方法不重新设计重建管线，而是直接在~Feature-3DGS~的双分支框架上构建语义层：外观分支照常渲染~RGB，特征场分支将高斯基元的语义信息编码为可渲染的特征向量。作者分别训练了~LSeg/CLIP、SAM1、SAM2~三种特征源，每种特征源独立监督一个特征场模型（各~7\,000~次迭代），以横向比较不同视觉基础模型在户外航拍场景中的语义表达能力。

\textbf{核心创新：语言引导的~3D~分割流水线。}这是本方法与锚点方法在功能层面的根本差异——Maikos et al.~的目标是快速生成高保真~3D~场景，而本方法的目标是让用户通过自然语言与该场景交互。具体流程为：给定文本提示（如「屋顶设备」「栏杆」），计算提示与~CLIP（Contrastive Language-Image Pre-training，对比语言-图像预训练）~嵌入特征场在新视角下渲染结果之间的余弦相似度热图；对该热图阈值化获得粗粒度~2D~分割掩膜；再以热图激活值最高的像素点作为提示点（prompt point）输入~SAM~或~SAM2，获得精细化的~2D~分割；最后在~3D~空间中通过余弦相似度聚类，将多个视角的~2D~分割结果聚合为对应目标物体的高斯基元子集，实现从「文本 → 热图 → 精细~2D~分割 → 3D~目标提取」的完整语义闭环。

\textbf{设计动机。}传统摄影测量（COLMAP、OpenMVS）可以输出精确几何，3DGS~可以输出照片级真实重建，但二者都缺乏场景级的语义可解释性——无法自动回答「这栋建筑上有没有裂缝」「楼梯栏杆在哪里」等问题。本方法在保持~3DGS~高质量渲染的同时，通过嵌入可查询的语言特征场，首次在~UAV~航拍场景中实现了灵活的语言驱动~3D~场景理解。

\textbf{实验结果。}在两个自采~UAV~基础设施数据集（Building: 223~张城市公园单层建筑，包含门窗和屋顶特征；Observatory: 253~张穹顶天文台，包含楼梯、栏杆和屋顶设备）上进行了验证。结果表明~CLIP-LSeg~特征场在捕获语义结构方面优于~SAM~系列，而~SAM2~在边界精度上优于~SAM1；整体验证了从航拍视角进行语言驱动~3D~场景理解的可行性。

\paragraph{Liang and Cheng：滑坡物理仿真}

与锚点方法~Maikos et al.~聚焦实时~3D~重建不同，Liang and Cheng（2026.05）~\cite{2605_10715}~的目标是将~3DGS~从静态场景表示扩展为物理仿真的动态场景表示——提出从~UAV~扫描到滑坡仿真的完整框架，将照片级真实场景捕获与基于~3DGS~的物理仿真相结合。锚点方法输出的是可实时渲染的~3D~场景，本方法则在此基础上追问：能否让~3DGS~模型直接参与滑坡等地质过程的物理仿真，在保持照片级真实感的同时可视化滑坡的动态演化？

\textbf{数据采集环节的差异。}Maikos et al.~使用~DJI Mini 4K~通过~RTMP~实时传输视频流，追求采集速度与实时性。本方法的数据采集策略则围绕~3D~重建质量设计：使用~DJI M3E~无人机，其~RTK~模块提供厘米级定位精度；拍摄遵循多距离（远/近距离）和多角度两个原则以确保几何的多尺度覆盖与稠密视角；此外，显式利用~UAV~记录的~GPS~元数据（WGS84~坐标系下的经纬度、高度和俯仰-偏航-横滚角），转换为~COLMAP~兼容的相机格式后作为位姿先验注入重建管线，从而将无人机高精度定位能力融入~3DGS~初始化。

\textbf{低各向异性正则化。}标准~3DGS~允许高斯椭球体具有极端纵横比，这在静态新视角合成中通常不影响渲染质量。然而，当高斯基元在物理仿真（如物质点法~MPM）中经历大幅平移和旋转时，高度各向异性的椭球体会产生拉长条纹或尖刺状渲染伪影。为此，本方法引入各向异性损失
\begin{equation}
L_{\text{aniso}} = \frac{1}{|P|}\sum_{p \in P}\bigl(\max\{(S_p^{\max})/(S_p^{\min}),\; r\} - r\bigr),
\end{equation}

其中~$S_p^{\max}$~和~$S_p^{\min}$~分别为第~$p$~个高斯的三个主尺度中的最大值和最小值，$r=3$~为允许的纵横比上界。该损失仅惩罚纵横比超过阈值的高斯，对近似各向同性的高斯不施加额外约束。最终损失函数为~$L = L_{\text{rec}} + \lambda_{\text{aniso}} L_{\text{aniso}}$。

\textbf{体积内部填充。}3DGS~重建仅捕获坡面的可见表面，但滑坡仿真需要完整的体积内部而非空心表面壳。本方法因此引入体积内部填充步骤：首先将重建高斯的中心视为点云，通过~Poisson~重建提取表面网格；然后在由目标仿真区域$(x_{\min}, x_{\max}, y_{\min}, y_{\max}, z_{\min})$定义的填充域内、以重建网格表面为顶部边界，均匀分布内部高斯中心以填充表面以下的土壤体积。每个内部高斯的尺寸根据其到最近表面点的距离~$d_p$~自适应设定：短轴~$s_p^{\text{short}} = d_p$，中轴~$s_p^{\text{mid}} = 1.5\,d_p$，长轴~$s_p^{\text{long}} = 2\,d_p$。颜色方面，为避免高斯基元在旋转过程中产生视角依赖的颜色不一致，内部填充高斯仅保留球谐函数的零阶项（基础颜色），丢弃所有高阶~SH~系数。

\textbf{MPM~滑坡仿真。}获得体积高斯模型后，本方法采用物质点法（Material Point Method, MPM）执行滑坡仿真。具体而言，为填充后的体积高斯基元分配材料属性、外力（如重力）和边界条件，由~MPM~求解器计算滑坡过程中每个时间步的粒子运动（变形、位移和流动）。由于场景直接由高斯基元表示，仿真产生的粒子运动自然传递到对应的高斯表示，从而可以在不同时间步渲染出滑坡演化的照片级真实可视化。仿真求解器基于~PhysGaussian~提供的~MPM~实现。

\subsubsection{高精度表面重建}

本路线介绍基于~3DGS/2DGS~的高精度航空表面重建方法。

\paragraph{AGS：航空表面重建}

Aerial Gaussian Splatting（AGS）（2024.09）~\cite{2409_00381} 与锚点方法~Maikos et al.~的端到端实时重建管线目标完全不同——锚点追求无人机飞行过程中的在线~3D~模型增量更新（渲染速度优先），而~AGS~的目标是利用~3DGS~实现大规模航空影像的离线高精度表面重建（几何精度优先），首次弥合~3DGS~在渲染质量上的优势与传统~MVS~方法在几何精度上的优势之间的鸿沟。

\textbf{原始~3DGS~用于航空表面重建的三重障碍。}（1）大规模航空场景（数百张千兆像素级图像、覆盖数万平方米）的庞大数据量导致~GPU~显存溢出；（2）3DGS~通过屏幕空间距离分配深度、基于~$\alpha$-混合~加权求和计算近似深度，无法精确确定射线与高斯的交点，因而无法获取可靠的法向量信息——而法向量是表面重建的核心输入；（3）原始~3DGS~仅使用图像级光度损失优化，高斯分布倾向于高保真渲染而牺牲表面几何精度。

\textbf{模块一：自适应航空场景分割。}AGS~借鉴~VastGaussian~的分块策略应对显存瓶颈，但针对航空场景的特殊几何做了关键扩展。先将场景按相机位置水平分成~$M$~块、垂直分成~$N$~块（训练前对场景执行~Manhattan~对齐，使~$y$~轴垂直于地面以方便分块），再扩展各块边界~20\%以减少拼接伪影。在此基础上，AGS~开发了锚点方法中未涉及的视角选择策略：将块内点云投影到所有图像上，若点落在图像中心~70\%~范围内则该图的投影分数加一，选取每区域投影分数最高的~$N$~张图像作为优化视角。这样做的目的是剔除对该块贡献微弱的远距视角、补充有效视角，避免因视角不足导致的局部重建错误。消融实验表明，关闭该策略后~$PAG_{0.6m}$~精度显著下降，部分区域的深度误差图出现严重错误，建筑物法线图出现透明和欠优化现象。

\textbf{模块二：射线-高斯交集。}这是~AGS~与锚点方法的核心差异之一。锚点方法沿用了标准~3DGS~的~$\alpha$-混合~深度近似，无需精确深度（其目标是渲染而非表面重建）；而~AGS~将相机中心~$o$~和射线方向~$r$~转换到高斯局部坐标系后，解析计算射线与高斯椭球的最大值交点深度
\begin{equation}
t_{\text{max}} = -\frac{o_{g}^{T}r_{g}}{r_{g}^{T}r_{g}}\,,
\end{equation}

替代标准~3DGS~的近似深度，并从交点处的切平面推导高斯法向量。获得精确深度和法向量后，AGS~引入~2DGS~的深度-法向一致性约束
\begin{equation}
L_n = \sum_i \omega_i (1 - n_i^T N)\,,
\end{equation}

其中~$\omega_i$~为混合权重，$n_i$~为高斯法向量，$N$~为从渲染深度图梯度计算的法向量。消融实验显示，引入射线-高斯交集后~$PAG_{0.6m}$~提升~19.60\%。

\textbf{模块三：多视角几何一致性约束。}标准~3DGS~（及锚点方法）仅在单视图内施加光度损失，容易过拟合且无法保证多视角几何一致性。AGS~对相邻视角~$V_r$~和~$V_n$~分别渲染深度图~$D_r$~和~$D_n$，将参考视角像素~$P$~通过深度和内外参投影到相邻视角得到~$P^*$，再反投影回参考视角得到~$P_r^*$，以两者坐标距离~$\|P - P_r^*\|$~作为几何一致性损失（距离超过阈值~1~像素的离群点置零以降低遮挡影响）。

\textbf{表面提取。}各数据块独立优化~50\,000~次迭代后（致密化从第~500~次开始、第~30\,000~次结束；几何约束从第~7\,000~次引入），去除扩展区域并合并，再通过截断符号距离函数（TSDF）从高斯场中提取最终网格表面——这一步骤沿用了~2D Gaussian Splatting~的表面重建管线。

\textbf{实验结果。}在~WHU-OMVS~数据集（268~张五镜头倾斜相机图像，550\,m~航高，覆盖~$850\times700$\,m$^2$，地面分辨率~10\,cm）上，AGS~在最严格指标~$PAG_{0.6m}$~上超越第二名~GOF~2.78\%、超越~COLMAP~11.26\%，在~MAE~和~RMSE~上同样达到最优。在弱纹理区域密集的天津数据集（342~张图像，200\,m~航高）上，AGS~在~0.6\,m~阈值下的精度比~2DGS~高~9.83\%、比~3DGS~高~26.10\%、比~GOF~高~2.32\%，达到与~COLMAP~相当的几何精度——这是~3DGS~类方法在航空表面重建中首次追平传统~MVS~的几何精度。在~Mill-19~和~UrbanScene3D~的大规模场景（1\,678--2\,582~张图像）中，AGS~不仅能重建建筑物整体，还能捕获窗户、空调外机、地面植被乃至路灯杆等精细结构。新视角合成方面，AGS~在~Building~和~Rubble~数据集上取得~PSNR、SSIM、LPIPS~三项指标的最优，在~LPIPS~上比~3DGS~提升约~46\%。

\textbf{锚点回顾。}Maikos et al.~建立了从无人机~RTMP~实时视频流到~3DGS~在线优化的全链路管线，验证了~3DGS~在航空场景中高速渲染和局部增量更新的可行性——但该管线（以及本章此前讨论的所有方法）均在单一光照条件下运行，未触及一个关键问题：当同一场景在不同时段（不同太阳高度角、阴影分布和曝光条件）下重复采集时，现有~3DGS~方法能否保持几何精度、渲染质量和材质估计的跨时段一致性？

\subsubsection{数据集与评测基准}

本路线介绍面向航拍~3D~重建评测的多时相光照基准和~UAV~智能体评测基准。

\paragraph{SkyLume：多时相光照基准}

SkyLume（2025.12）~\cite{2512_14200}~不提出新的重建或渲染算法，而是构建一个面向变光照~3D~重建评估的大规模~UAV~基准数据集——其核心任务是将"多时相采集"从工程难题转化为可量化评估的实验条件。具体而言，SkyLume~的贡献在于三方面：（1）提供严格视角一致的多时段数据，使时段间差异可归因于光照而非几何；（2）以~LiDAR~点云为锚点实现多时相亚像素级对齐；（3）引入时间一致性系数（TCC）作为衡量逆向渲染跨时段稳定性的新指标。

\textbf{数据采集设计。}SkyLume~包含~10~个城市区域超过~10~万张高分辨率多视角图像（CHCNAV C30~倾斜相机，130\,MP~五镜头同步采集，前向重叠~80\%、侧向重叠~60\%，120\,m~航高，8\,m/s~航速），每个区域在三个时段（早晨~07:00--09:00、正午~11:00--13:00、下午~16:00--18:00）以完全相同的航迹重复采集。这一设计的直觉含义是：若三个时段的相机视角完全一致，则渲染结果之间的差异只能来自光照——这正是将光照鲁棒性转化为可测量属性的关键。LiDAR~数据（DJI Zenmuse L2，水平~5\,cm、垂直~4\,cm~精度）在单独一天采集，因为~LiDAR~作为主动测量模态不受光照变化影响，可充当跨时段对齐的几何锚点。

\textbf{多时相对齐管线。}数据集的核心技术贡献在于多时相精确对齐——这一管线分为四步串联：（1）将~LiDAR~数据在~WGS84/UTM~坐标系中地理参考化；（2）通过~RealityScan~注册~LiDAR~点云并生成固定相机位姿作为对齐锚点，渲染伪~RGB LiDAR~视图；（3）将三个时段的~RGB~图像强制与~LiDAR~位姿对齐；（4）执行联合~SfM~（RTK~位姿先验作为软约束，初始化相机位姿后在全局光束法平差中精化），对稀疏特征区域辅以少量地面控制点（GCP）消除残余漂移。最终对齐精度：中位重投影误差仅~0.70~像素，相机位置不确定性在毫米级。在弱纹理和阴影区域，LiDAR~提供的度量支架（metric scaffold）有效正则化深度估计，使融合~LiDAR+RGB~的稠密重建显著优于纯~RGB~重建。此外，数据集对河流和池塘等水体区域进行专门处理——通过~LiDAR~在岸边采样四点计算平均高程（样本散布~$<$2\,cm），将法向量正则化为近水平方向后重新网格化，确保水面几何与周围地形的连续性。

\textbf{输出数据与评估协议。}数据集以~COLMAP~格式输出每个时段的相机外参/内参和稀疏点云，并提供两种模态的逐帧深度和法向量图：（a）网格深度/法向量（从融合网格投影至各相机视角，分辨率与原始图像一致）；（b）LiDAR~深度（LiDAR~点云经占用滤波后重投影至相机帧，虽更稀疏但度量可靠性更高）。数据集还提供逐帧太阳几何信息（通过~NOAA~太阳几何计算器从时间戳和地理位置计算太阳高度角和方位角），以支持去阴影、逆向渲染和重光照研究。

\textbf{TCC~指标设计。}SkyLume~引入时间一致性系数（Temporal Consistency Coefficient, TCC）作为评估逆向渲染跨时段稳定性的核心指标。其直觉含义是：对同一场景在不同光照下训练的模型，若逆向渲染成功分离了光照与材质，则从同一测试视角渲染的~albedo~图在三个时段间应完全一致。具体地，TCC~从全局轨迹中选取~$K$~个空间均匀分布的固定测试视角，在每个视角渲染三个时段的~albedo~图，以四个子指标的加权组合评估跨时段一致性——$\text{TCC}^{(k)}_\text{comb} = 0.2\,\text{TCC}_\text{MAE} + 0.2\,\text{TCC}_\text{RMSE} + 0.2\,\text{TCC}_\text{SSIM} + 0.4\,\text{TCC}_\text{LPIPS}$（感知分支~LPIPS~权重最高为~0.4，反映对视觉可信度的重视），最终~TCC~得分为所有测试视角的算术平均。这一设计的核心创新在于：TCC~奖励的是跨时段稳定性而非单时段保真度——即使某个时段的~albedo~估计再精细，若其他时段出现漂移，TCC~分数仍会显著下降。

\textbf{基准实验发现。}在~6~个中小规模场景上对三类方法（逆向渲染、几何提取和新视角合成）的基准评估揭示了三个关键挑战：（1）同一区域内光照变化导致渲染质量大幅波动，弱纹理立面、玻璃边界和临水表面尤为严重；（2）几何常被光照效应"腐蚀"——投射阴影和移动半影被过拟合为实体结构，偏置了深度和法向量估计，引入虚假几何体；（3）城市级逆向渲染仍然脆弱——估计的~albedo~保留了时段相关的阴影和曝光变化，未能实现跨时段一致性。这些发现为光照鲁棒的大规模~3D~重建指明了方向：必须显式分离光照与材质属性，并引入处理动态光照和确保跨时段一致性的专门技术。

与~SkyLume~聚焦光照评估不同，下面转向智能体行为的评估。

\paragraph{HUGE-Bench：UAV智能体评测}

HUGE-Bench（2026.03）~\cite{2603_19822} 将评估目标从"重建质量"前移至"智能体行为"——其核心问题是：当~UAV~重建已经可行（如上文~Maikos et al.~的实时~3DGS~管线所示），智能体能否仅凭简短自然语言指令（如"检查左侧建筑"）自主分解并安全执行多阶段飞行任务？

\textbf{锚点回顾与增量差异。}上文~Maikos et al.~解决了"无人机边飞边重建"的问题——用~3DGS~在线优化实现实时场景重建。HUGE-Bench~在此基础上向前迈进一步：它不再评估~3DGS~本身的重建质量（PSNR/SSIM），而是将~3DGS~重建产物用作仿真环境的视觉层，测试~UAV~智能体在高级语义命令下的行为执行能力。这一转变的关键在于，HUGE-Bench~要求~3DGS~不仅渲染得"像"，还要支持物理安全的碰撞检测——因此引入了与~3DGS~空间对齐的三角网格（triangle mesh），构成双模态数字孪生。

\textbf{对齐~3DGS-Mesh~双模态数字孪生。}这是~HUGE-Bench~的核心基础设施，也是与纯~3DGS~重建管线的关键区别：对每个场景同时重建~3DGS~模型和三角网格，两者在空间上精确对齐。3DGS~负责高保真渲染，为~VLA~策略提供真实感知输入（RGB~图像）；Mesh~负责碰撞检测几何和基于光线投射的深度查询。这种对齐表示使仿真能同时兼顾视觉真实性（3DGS~渲染接近真实照片）和物理安全性（Mesh~提供可碰撞查询的几何体）。数据通过~DJI M400 UAV（配备~Zenmuse L2~载荷）在四个室外场景（办公建筑、密集城市街区、沼泽农田、施工道路，总面积约~6.45~km$^2$）采集，飞行高度约~85\,m，记录~GNSS~定位、UAV~位姿、相机位姿和高分辨率~RGB~图像（$5280\times3956$），并补充近地低空视角以提升重建质量。

\textbf{高级~VLA~任务定义。}与现有~UAV VLN~基准使用冗长逐步路线描述不同，HUGE-Bench~的指令风格短促且含歧义——更接近真实操控者的指令习惯。基准定义了八种高级~UAV~任务：目标着陆（Landing）、道路巡检（Inspection-R）、自适应建筑巡检（Inspection-B）、区域测绘（Mapping）、不同高度轨道飞行（Orbit-H）、不同半径轨道飞行（Orbit-R）、多圈螺旋下降（Spiral Down）和障碍物规避区域穿越（Traversal）。每条短指令隐含多步子任务序列——例如"检查左侧建筑"要求~UAV：识别目标建筑$\to$接近$\to$下降至巡检高度$\to$保持安全距离环绕飞行$\to$返回初始位置。这些任务共同考察目标识别、方向接地（directional grounding）、深度与高度感知以及多阶段轨迹完成能力。数据集共包含~5\,330~条训练轨迹、593~条测试可见轨迹和~294~条测试不可见轨迹（2.56M~米总轨迹长度），后者通过保留未见地标和~LLM~指令复述（paraphrase）评估模型对新场景和新语言的泛化能力。

\textbf{Real-to-Sim~管线与标注生成。}HUGE-Bench~通过一条从真实到仿真的管线生成轨迹数据：(1) 真实场景采集；(2) 对齐~3DGS-Mesh~数字孪生重建；(3) 标注生成——在~60\,m~高度部署俯视相机网格，渲染~3DGS~获取全局~2D~地图视图并渲染~Mesh~获取对齐深度，再由~LLM~在渲染地图上定位地标并输出归一化图像坐标，经人工审核后反投影为~3D~地标位置；对于建筑、道路等语义扩展区域还提供人工标注掩码；(4) 在~Isaac Sim~中基于~3DGS-Mesh~数字孪生进行仿真轨迹采集——使用~GPU~加速渲染和~PhysX~刚体仿真，记录同步的多模态数据流（RGB、深度、位姿、飞行状态、碰撞信号）。一个特别的设计细节是空间方向描述：当场景中存在多个视觉相似的候选目标时，仅用"建筑""道路"等通用名词不足以唯一指定目标——因此系统在轨迹采集前选择目标地标，在其附近随机采样初始视角，使用目标掩码提示~LLM~生成类人空间指代表达（如"左上方的建筑""画面右侧沿线的道路"），使指令结合方向短语和通用类别实现更精确的语言接地。

\textbf{面向过程的评估体系。}传统~VLN~基准主要关注终端成功（是否到达目标），但~UAV~高级任务的失败模式更复杂——策略可能到达终点却跳过了必需的巡检程序，偏离了预期过程，或在途中发生碰撞。HUGE-Bench~因此提出三维度评估：(1)~\textbf{过程保真度}——轨迹覆盖率~TCR（Trajectory Coverage Rate）衡量预测轨迹对真值过程的覆盖程度，而非仅看终点，对巡检、测绘、轨道飞行等过程关键型任务报告~TCR@$K$（$K=1,2,5$\,m~为距离容差阈值）；(2)~\textbf{终端准确性}——成功率~SR（Success Rate）用于目标着陆等以到达终点为核心的任务；(3)~\textbf{安全性与效率}——碰撞率~CR（Collision Rate）和碰撞感知~SPL（CSPL，Collision-aware Success weighted by Path Length）联合衡量目标到达、无碰撞执行和路径效率。基准测试了~OpenVLA、$\pi_0$、$\pi_{0.5}$~和~FastVLM~四个代表性~VLA/VLM~模型，结果显示现有模型在短指令下的过程完成和安全执行方面仍存在显著差距——即使最强的~$\pi_{0.5}$~在过程关键型任务上也远未饱和，定位~HUGE-Bench~为高级~UAV VLA~的诊断性和安全性测试平台。

\begin{takeaway}
航空影像重建将 3DGS 应用于中等距离高分辨率观测条件：对齐 3DGS-Mesh 双模态数字孪生同时兼顾视觉真实性和物理安全性，为无人机视觉语言导航等高级任务提供了仿真评测基础设施。
\end{takeaway}

%=====================================================================

\subsection{近期方法脉络：UAV facade digitization with Gaussian Splatting}

这一组近期工作延续 Gaussian Splatting，主要把改动集中在应用等环节。下面只保留相对既有路线的变化，便于和前文方法直接对照。

\paragraph{RefGlass-GS}
RefGlass-GS（2026.06）~\cite{2606_28826} 延续 Gaussian Splatting 在“UAV facade digitization with Gaussian Splatting”上的基本设定，主要改变是：UAV 视角规划、玻璃面板分割与 Gaussian Splatting 融合。3DGS 的显式高斯表示和可微渲染能力保持为基础。
其设计动机是：针对反光玻璃导致的几何畸变、视角相关纹理和对象语义增强困难，提出端到端融合框架。
具体而言，包括基于结构规则的玻璃面板分割、最大覆盖的 UAV 视角规划，以及用于真实感渲染和交互展示的 Gaussian Splatting 表示。
适用边界是：源码章节未抽取，分类依据摘要；反射内容和透明材质的物理一致性可能仍难以保证。
\section{文化遗产与测绘}
\label{sec:20_3}

文化遗产数字化保护和测绘制图是~3DGS~在遥感中日益重要的应用方向。核心需求包括大规模场景的高保真~3D~重建、真实数字正射影像图（TDOM，True Digital Orthophoto Map）生成和交互式虚拟漫游。

\subsection{文化遗产与测绘方法}
\label{ss:20_4}

本节讨论~2~条技术路线：（1）文化遗产数据集与数字化；（2）数字正射影像图与测绘。

\subsubsection{文化遗产数据集与数字化}

本路线介绍文化遗产场景的大规模~3DGS~重建数据集、轻量化重建管线和旧照片档案利用方法。

大规模文化遗产场景的~3DGS~重建面临数据规模、采集条件和平台适配等多重挑战。本节从数据集基准建立、轻量化重建管线和旧照片档案利用~3~个方面展开，涵盖~4~种代表性方法。

CULTURE3D（2025.01）~\cite{2501_06927} 发布了包含~41006~张无人机高分辨率航拍图像的超大型细粒度数据集，覆盖~20~个全球文化地标（如剑桥大学、金字塔、故宫），总点数达~100~亿，由此建立了大规模高斯~Splatting~评估的全新基准。此前，大尺度~3D~重建基准面临三重瓶颈：MatrixCity~等合成数据集与真实场景存在显著的域迁移差距；GauU-Scene~虽提供真实航拍数据但场景同质化严重（仅户外环境，约~46000~张图像、6.28~亿点），不足以暴露大规模重建的内存瓶颈；文化遗产场景在已有数据集中被严重忽视。CULTURE3D~正是针对这些不足而设计——以真实世界的大体量（100~亿点级）、高精细度（~48MP~图像分辨率~8064$\times$6048）和场景多样性（室内外混合，涵盖佩特拉古城、龙门石窟、卢浮宫等~20~处地标）来"压力测试"当前方法。

在实现上，CULTURE3D~建立了从数据采集到基准评估的完整管线。采集端使用~DJI Mini 3~无人机以约~15~度角度间隔的受控轨道和网格飞行模式采集图像；重建端先以~COLMAP~进行~SfM~稀疏重建，再以~Reality Capture~完成~MVS~稠密重建生成纹理网格和相机参数；评估端在~8$\times$NVIDIA RTX A6000（48~GB/GPU）上统一训练~30000~次迭代，对~7~种主流方法使用相同输入进行公平对比。基准结果系统性揭示了当前方法的可扩展性瓶颈：标准~3DGS~训练最快（$<$1~小时）但质量最低（SSIM~仅~0.2861--0.6268）；SuGaR~和~Wild Gaussian~在~Cambridge~场景直接~OOM；GOF~在~Pyramid~场景彻底~FAIL。唯一在全部测试场景上无~OOM/FAIL~的方法是~City GS（PSNR~14.37--21.86~dB），其分层高斯组织策略有效管理了大量基元的内存开销。

与锚点~CULTURE3D~以专业无人机和工业级重建管线建立大规模评估基准的思路形成互补，Gaussian Heritage（2024.09）~\cite{2409_19039} 从应用的另一端切入。它回答的问题是：资源有限的文化遗产机构能否仅凭一部智能手机和~RGB~图像，即可完成~3D~场景重建并将单个文物从场景中自动提取出来？这一设计动机源于对文化遗产数字化现实的观察：当前艺术数字化方法缓慢、昂贵，往往依赖专业技术人员和专用设备，严重限制了这些技术在中小机构中的普及。因此，本方法刻意将数据需求降至最低——仅需标准智能手机拍摄的~RGB~图像，无需深度传感器、无需手动标注、无需专业采集设备。

为实现"重建$+$分割"的双重目标，本方法在标准~3DGS~之上引入分割感知能力：为每个高斯基元增加~16~维特征向量编码分割信息，该特征通过~ViT-H SAM~自动生成的~2D~分割掩码进行监督训练，无需人工标注。训练过程联合优化三个损失——标准~3DGS~渲染损失（几何与外观）、对比聚类损失（将渲染特征与~2D~分割掩码对齐）和空间相似性正则化（鼓励空间相邻高斯具有一致的分割特征、空间远离高斯具有不同特征）。分割阶段支持两种模式：实例分割通过在任意视角渲染分割特征后以余弦相似度阈值（$t=0.7$）聚类像素；语义分割则接收文本提示，通过~Grounding DINO~生成边界框并由~SAM~精炼为掩码。最终~3D~提取通过渲染初始视图、用~GroundingDINO+SAM~生成目标物体掩码、基于余弦相似度分割~3D~高斯子集并施加凸包算法改善分割连贯性。在~LERF-Mask~和~3D-OVS~两个标准~3D~场景分割基准上，本方法在~mIoU~和~mBIoU~指标上均优于~LeRF、Gaussian Grouping~和~LangSplat~等基线方法。整个管线以~Docker~容器封装部署，从~RGB~图像输入到可交互的~3D~文物模型输出实现全自动化。

与~CULTURE3D~聚焦大规模基准评估的定位不同，Nevis Digital Twin（2026.03）~\cite{2603_20560}~将研究视角从"评估"转向"应用"——为加勒比海地区尼维斯岛的濒危遗产建筑建立多模态数字孪生并融合到沉浸式~VR~环境中。如果说~CULTURE3D~以受控无人机飞行定义了理想化的采集条件，Nevis Digital Twin~则代表了完全不同的操作约束：在热带岛屿环境下，仅凭消费级设备实现遗产场景的快速数字化。本方法的核心贡献在于数据采集和虚拟重建的多模态融合策略。

数据采集方面，设计了三级互补方案：（1）Timelapse~模式——使用~Insta360 X4 360\textdegree~相机以~5.7K@2fps~采集，数据量较小但低帧率可能导致快速移动时出现覆盖间隙；（2）Video~模式——5.7K@30fps~视频采集，高帧率确保~SfM~对齐所需的冗余帧以消除运动模糊，成为最终工作流的首选模式；（3）精细扫描——对单体建筑和高价值文物使用~iPad Pro LiDAR~进行轨道式扫描（眼平、高角度、低角度三圈环绕），利用~PolyCam~在~iPad~端生成亚毫米精度网格。所有~360\textdegree~数据以~2~米间距~zig-zag~路径步行采集，相机固定于隐形自拍杆顶端约~1~米过头高度，全向视野从根本上消除了传统摄影测量中精心瞄准以确保~60--80\%~重叠的需求。

Scan-to-BIM~流程方面，原始双镜头素材先经~Insta360 Studio~拼接为等距柱状视频，再按步行速度计算提取率以保证至少~70\%~前向重叠，提取帧上传至~CupixVista~云端平台通过~SLAM~算法对齐生成~3D~网格和可导航的~Street View~风格导览；高价值文物则通过~PolyCam LiDAR~扫描的精细网格嵌入到整体场景模型中，形成多尺度数字孪生环境。虚拟重建方面，同时探索了网格重建和~3D Gaussian Splatting~两种技术路线：3DGS~管线基于~Nerfstudio~框架的~Splatfacto~模型实现，使用~COLMAP~进行~SfM~相机位姿估计生成稀疏点云作为高斯初始化，底层依赖~gsplat~的优化~CUDA~内核进行可微分渲染和快速收敛；3DGS~在照片级真实感和复杂场景处理方面展现出独特优势，而网格在几何精度和~Scan-to-BIM~工作流兼容性方面更为成熟。

VR~部署方面，训练好的~3DGS~模型导出为标准~.PLY~格式，在~Unity~引擎（2022.3.62f2）中基于开源~UnityGaussianSplatting~框架导入并转换为可调压缩级别的优化资产以降低显存占用，通过~OpenXR~和计算着色器在~Meta Quest 3~头显（单眼~$2064\times2208$~分辨率）上实现沉浸式~VR~可视化，开发测试在~AMD Ryzen 9 5900HX / RTX 3070（8\,GB VRAM）笔记本平台上完成——整体为濒危遗产的远程访问和保护决策提供了从数据采集到~VR~交付的完整工作流。

与~CULTURE3D~使用~DJI Mini 3~无人机以受控轨道和网格飞行模式系统采集不同，Shih（2025.03）~\cite{2503_10031} 面临的挑战更为极端——利用~17~年前的旧照片重建台湾云林 Gongfan Palace（建于~1685~年、2006~年被指定为县级历史遗址）的装饰艺术细节。这些照片完全不具备~CULTURE3D~那样的受控采集条件：约~1500~张影像由~SONY DSC-R1（3888$\times$2592~像素）和~Panasonic DMC-FX100（4000$\times$3000~像素）消费级相机拍摄，采用径向和向心式构图记录工作过程和环境观察，既无位姿信息也无系统化的拍摄规划；照片中还大量包含用于激光扫描配准的条纹测量标记，对~3D~重建造成额外干扰（多数标记需自动或手动移除）。这一场景代表了文化遗产数字化中的一个普遍现实：许多已损毁或改造的历史建筑仅存业余照片档案，缺乏专业的~3D~采集数据——该庙宇早在~2008~年就曾使用~Leica HDS 3000~激光扫描仪进行过专业采集（完成~612~次扫描、7450万点），但设备笨重、空间受限且成本高昂，而此次工作则探索能否以旧照片替代专业设备实现可用的~3D~重建。

本方法的核心贡献在于多平台~AI~辅助建模的系统性对比评估。在重建平台上，三款工具覆盖了不同的技术路线：Postshot（桌面端，基于~Nvidia RTX 2060~本地训练~3DGS，解决了数据隐私和主权问题）、KIRI Engine（云端计算，支持~3DGS~和~NeRF~两种输出模式，最低~20~张图像即可启动）和~Zephyr（传统摄影测量管线，生成多边形/点模型作为参考基准）。重建目标涵盖庙宇的五大类装饰构件——石雕（龙柱、石狮、壁面装饰，含镂空、浮雕和圆雕三种技法）、庙宇结构（三川殿、屋顶、室内）、木雕（雀替、垂花、吊筒）、泥塑（屋脊下方壁面）和剪粘工艺——通过~6~项评估标准进行跨平台质量检验：可用图像充足性、细部层次一致性、图像分割准确性、多边形输出选项、程序可访问性和种子值可复现性。定量评估结果揭示了不同平台间的显著差异：KIRI Engine~与~Zephyr~配对比较的标准差仅~2.93~mm（最大距离~17.80~mm，平均距离~3.23~mm），而~Postshot~与~Zephyr~配对比较的标准差则高达~21.69~mm（最大距离~140.25~mm，平均距离~19.13~mm）——Postshot~导出的~PLY~格式高斯呈不均匀分布的飞溅状，与传统多边形模型的几何对齐精度明显较低。总体~3D~输出成功率约~80\%，但满意成功率仅约~50\%，反映了非结构化旧照片数据对当前~3D~重建技术的真实压力。

此外，该工作还将~AI~辅助建模延伸到下游应用：整合~Stable Diffusion~进行~AI~辅助图像生成与增强，通过~Postshot~直接从图像生成飞穿动画（无需传统多边形建模中间步骤），并以~ComeTrue T10~彩色~3D~打印机（1200$\times$556~dpi，喷墨染料逐层打印石膏粉末）验证物理输出质量——旧照片重建的细节通过提高模型分辨率得到改善。这项工作展示了~AI~辅助工具如何将非结构化的历史照片档案转化为~3D~重建资源，其核心价值在于提出了一种将遗产文档与~AI~技术连接的连续保护工作流程——从旧照片到~3DGS~模型，再到游戏引擎（Unreal Engine）、多媒体（After Effect）和~AR~交互（Augment）的多样化输出接口，形成了传统~3D~扫描管线无法覆盖的轻量化替代路径。

\subsubsection{数字正射影像图与测绘}

本路线介绍基于~3DGS/2DGS~正交投影的~TDOM~生成、交互式测量和建筑网格自动提取方法。

真实数字正射影像图（TDOM）生成和精确测绘是~3DGS~在遥感测绘中的核心产品方向。与自然场景不同，航拍测绘需要将~3DGS~的透视投影转换为正交投影，并解决大规模场景分块处理和增量更新等问题。本节介绍~5~种代表性方法，涵盖~3DGS/2DGS~正交投影渲染、多射线交互式测量和建筑网格自动提取等方向。

TOrtho-Gaussian（2024.11）~\cite{2411_19594} 首次提出了通过正交投影~3DGS~高斯核生成~TDOM~的方法。与传统摄影测量管线需要先生成数字表面模型（DSM）再进行遮挡检测和正射纠正不同，该方法直接将~3DGS~的~3D~高斯基元通过正交投影映射到地面平面，利用高斯的不透明度加权自然处理遮挡关系——距离相机更近的高斯具有更高的渲染权重，自动遮挡后方高斯。采用分而治之策略处理大规模区域。与锚点~TOrtho-Gaussian~将~3DGS~高斯基元直接用于正交投影渲染~TDOM~的产品输出不同，Deng and Qin（2026.03）~\cite{2603_24716}~将~3DGS~的角色从渲染引擎转变为高保真视觉参考——利用其无缝新视角合成能力作为操作者瞄准测量射线的交互界面，而非生成任何正射影像产品。这一设计动机在于：传统立体测量需要专业摄影测量工作站和密集匹配生成的网格模型，而网格表面的粗糙采样在薄结构和锐角处无法精确表示尖锐几何特征，导致直接在网格上的点选测量精度受限；3DGS~渲染的照级真实感使操作者能够跨视角准确识别同一语义点（"一致点"），绕过了底层网格的几何不准确性。具体算法流程为循环式多射线空间前方交会：操作者在~3DGS~场景中自由导航到第$i$个视角$V_i$，点击像素$p_i$标记一致点，渲染引擎返回当前相机位姿（相机中心$C_i = (X_i, Y_i, Z_i)^T$和从$C_i$穿过$p_i$的方向向量$d_i = (u_i, v_i, w_i)^T$），形成射线$R_i = (C_i, d_i)$；重复此过程收集$N \geq 2$条射线后，构建超定线性最小二乘系统。每条射线的几何约束来自叉积条件$(P - C_i) \times d_i = 0$，展开后得到三个方程但仅有两个线性独立，选取前两个整理为$2 \times 3$矩阵$A_i$和$2 \times 1$向量$\mathbf{b}_i$：

\begin{equation*}
A_i = \begin{pmatrix} 0 & w_i & -v_i \\ -w_i & 0 & u_i \end{pmatrix}, \quad
\mathbf{b}_i = \begin{pmatrix} w_i Y_i - v_i Z_i \\ -w_i X_i + u_i Z_i \end{pmatrix}
\end{equation*}

$N$~条射线堆叠为$2N \times 3$矩阵$A$和$2N \times 1$向量$\mathbf{b}$，最小二乘解为$\mathbf{x} = (A^TA)^{-1}A^T\mathbf{b}$。精度评估框架基于残差$\mathbf{v} = A\mathbf{x} - \mathbf{b}$，以系统自由度$r = 2N - 3$计算后验单位权方差$\hat{\sigma}_0^2 = \mathbf{v}^T\mathbf{v}/r$（其平方根~RMS~即射线到计算点的平均几何距离，作为每个测量点的核心不确定性指标），并推导协方差矩阵$Q_x = \hat{\sigma}_0^2(A^TA)^{-1}$，其对角元素提供各坐标分量的方差$(\sigma_x^2, \sigma_y^2, \sigma_z^2)$。值得注意的是，论文指出虽然可以生成任意数量的合成观测视角，但这些额外采样不会增加超过原始图像所提供的有效自由度——增加采样视角数主要提升数值稳定性而非引入新的独立信息。整个方法被实现为基于~Cesium.js~平台的~Web~端概念验证应用，操作者可在标准浏览器中完成全部测量流程。在薄结构和锐角等挑战性几何上，该方法显著优于网格直接测量（RMSE~达~1--2 cm~级别），因为~3DGS~渲染的高保真度使操作者能跨视角准确识别一致点，而网格表面的粗糙采样往往无法精确表示尖锐几何特征。

与锚点~TOrtho-Gaussian~的离线批量~TDOM~生成不同，A-TDOM（2025.09）~\cite{2509_12759} 实现了近实时增量式~TDOM~生成——在每张新图像到达时即刻更新~TDOM~输出，其设计动机在于将~TDOM~从后处理产品转变为可与数据采集同步进行的在线服务。

A-TDOM~的管线分为三个阶段：阶段~I~等待~$N_0$~张初始图像完成~On-the-Fly SfM~后初始化~3DGS~场，为防止仅用~2--3~张初始图像导致过拟合，采用紧约束的训练迭代数量；阶段~II~为增量更新核心——每张新图像到达时执行四步操作：（1）通过~On-the-Fly SfM~注册新图像位姿并精炼已注册图像位姿，产生更新的稀疏点云；（2）将更新后的稀疏点云重投影到新图像上进行~2D Delaunay~三角化确定关键区域掩码——这一设计至关重要，因为标准~3DGS~优化会试图重建仅被单视角或双视角覆盖的区域（缺乏多视角约束），导致不稳定、不一致的几何和外观，而~Delaunay~掩码利用~SfM~过程中~bundle adjustment~仅保留高质量点这一特性，确保只优化有多视角三角化充分覆盖的区域；（3）基于~Laplacian of Gaussian（LoG）梯度差异图识别需要新高斯基元的区域——对当前~3DGS~渲染图像和输入裁剪图像分别计算~LoG~梯度图，两图相减得到梯度差异，超过阈值~$G_m$~的像素标记为需要新增高斯的区域~$R_a$——然后在~Delaunay~三角形内均匀采样~2D~点，通过重心坐标插值投影到~3D~空间初始化新高斯的位置$P_S = \lambda_1 P_1 + \lambda_2 P_2 + \lambda_3 P_3$和颜色$C_S = \lambda_1 C_1 + \lambda_2 C_2 + \lambda_3 C_3$（替代标准~3DGS~的梯度驱动~split/clone~致密化控制，提供更高的可控性并避免实时应用中的过度高斯累积）；（4）增量~3DGS~优化后通过正交~splatting~生成更新~TDOM。

在自适应优化策略方面，A-TDOM~对标准~3DGS~训练进行了三项关键修改：迭代分配——最新输入图像分配一半训练迭代~$T_I/2$~（因最新图像通常包含最大比例的未见过区域），剩余~$T_I/2$~均分给已注册图像以保证持续优化；学习率按图像独立分配和衰减——渲染质量高的图像分配较低学习率以稳定精炼，渲染质量差的图像使用较高学习率以加速收敛；最终精炼阶段禁用高斯致密化策略，仅优化现有高斯以确保不引入冗余基元。阶段~III~在所有图像采集完成后进行轻量级精炼——训练迭代均分给所有图像，学习率仍按图像独立分配（与阶段~II~一致），使渲染质量较低的区域获得更积极的优化。在正交投影实现上，A-TDOM~分析了~3DGS~透视投影矩阵~$P$~的三个组成变换：视锥-长方体变换~$M_{\text{persp2ortho}}$、长方体-规范立方体变换~$M_{\text{ortho}}$~和规范立方体-NDC~变换~$M_{\text{hand}}$，通过移除非仿射的~$M_{\text{persp2ortho}}$~项得到正交投影矩阵~$P_0$，由于正交投影是仿射变换，其~Jacobian~可直接从~$P_0$~解析获得~$J_0 = \text{diag}(2/(r-l), 2/(t-b), 0)$，避免了透视投影中~Jacobian~的非线性近似。

渲染分辨率由地面采样距离（GSD）确定：给定场景包围盒尺寸和目标~GSD，输出图像尺寸计算为~$W = (X_{\max} - X_{\min})/\text{GSD}_x$~和~$H = (Y_{\max} - Y_{\min})/\text{GSD}_y$。

与锚点~TOrtho-Gaussian~基于~3DGS~的正交投影不同，Wang et al.（2025.03）~\cite{2503_19703} 基于~2DGS~提出免~DSM~和遮挡检测的~TDOM~生成方法——这一选择源于~2D Gaussian Surfel~的平面先验比~3D~椭球体更适合表面重建任务，天然避免了~3DGS~中因视角变化导致的深度不一致问题（3D~椭球体在不同视角下的形状变化导致表面解释不一致）。在基元表示上，2DGS~为每个高斯构建局部坐标系：任一点$P(u,v)$通过高斯中心$\mu$、两个切向量$(t_u, t_v)$和缩放因子$(s_u, s_v)$变换到世界坐标，其中法向量$t_w = t_u \times t_v$定义局部平面方向——这种显式平面参数化使~2DGS~在表面建模中比~3D~椭球体更精确。本方法对~2DGS~正交投影的关键修改在于：将~2DGS~原始的透视投影矩阵$M_{\text{persp}}$替换为正交投影矩阵$M_{\text{ortho}}$——透视投影中所有射线从光心出发经视锥映射到~NDC~空间，而正交投影中所有射线平行于光轴，将长方体视域直接映射到~NDC~空间。

由于~2DGS~的~2D~高斯基元是平面盘状结构（而非~3DGS~的椭球体），正交投影下每个点的变换仅需$M_{\text{ortho}}WH$（投影矩阵$\times$视图矩阵$\times$局部-世界变换），无需~3DGS~正交~splatting~中对协方差矩阵的额外仿射变换步骤，简化了投影管线。在正交深度图生成方面，本方法通过体渲染公式$D = \sum_i \alpha_i T_i d_i$计算每个像素的期望终止深度，其中$\alpha_i$为射线-高斯交点$x$处的高斯值$G_i(x)$（近似不透明度），$T_i = \prod_{j=0}^{i-1}(1-\alpha_j)$为透射率，$d_i$在正交模型中等于射线-高斯交点在视图空间的~z~轴坐标（因为所有射线平行于光轴，射线起点为经过相机光心且垂直于射线的平面上的交点）。针对~2D~高斯基元在视角接近平行于高斯平面时退化为线段的问题（此时交点处的高斯值$G_i(x)$过小导致深度估计不可靠），本方法引入退化检测机制：在图像平面上以投影高斯中心为均值构造标准~2D~高斯$g_i^{2D}(x)$，当$g_i^{2D}(x) > G_i(x)$时判定当前视角下该高斯退化，改用高斯中心的视图空间~z~坐标作为深度$d_i$——这确保了即使在极端视角下深度图仍然可靠。

大规模场景处理采用与~TOrtho-Gaussian~类似的分而治之策略，具体包含三步：（1）基于相机位置的场景分块——将相机中心投影到地面平面，划分为$m \times n$个网格单元并确保每个单元内相机数量相等；（2）基于位置的点云选择——对每个单元边界按预设比例扩展后，将点云投影到地面平面并选择扩展边界内的点作为当前分区的输入数据；（3）基于可见性的相机选择——定义第$j$个单元对第$i$个视角的可见性为$\phi_{ij}/\Phi_i$（$\Phi_i$为第$i$个视角的图像面积，$\phi_{ij}$为第$j$个单元包围盒在第$i$个视角中的投影面积），当可见性超过阈值时该相机被选入当前单元，其覆盖的点也被纳入。训练前先通过~COLMAP~稀疏重建获取初始点云和相机位姿，经~Manhattan~对齐将场景包围盒与坐标轴对齐，每个场景分为~4~个子区域在~4~块~RTX 3090~上各自训练~60,000~次迭代，平均训练时间约~7~小时。

训练后按分块包围盒过滤高斯中心位置并进行拼接去重，渲染时将全图分为$2 \times 2$切片逐块渲染，同时为每个像素计算正交深度，支持批量光栅化生成与输入图像分辨率相当的正射全景。在定量评估中，本方法在~NPU DroneMap Dataset~（最大场景覆盖~1.598 km$^2$）上通过~GCP~距离误差评估几何精度——在与~GCP~对齐后，大部分误差控制在~1.0 m~以内；在~ENU~坐标系下的~RMSE~评估中，本方法达到~1.928 m，略优于锚点~TOrtho-Gaussian~的~2.003 m，且优于~Metashape（2.315 m）、Map2DFusion（3.106 m）和~Pix4Dmapper（3.798 m），仅次于~ContextCapture（0.701 m）。在图像重建质量方面，本方法的~SSIM~与~TOrtho-Gaussian~相当，PSNR~和~LPIPS~略低——这在以俯视航拍为主的设置下（缺乏多角度约束）属于预期范围，本方法的核心关注点在于~2DGS~表示平面表面的能力及其对正射影像精度的提升。

在建筑物和薄结构（电力塔、空调护栏）的定性对比中，本方法在边缘锐度、遮挡关系恢复和细节保持方面优于~Map2DFusion、ContextCapture、Metashape~和~Pix4Dmapper~四款商业软件以及~TOrtho-Gaussian——后者虽保持整体结构但缺乏高频细节，降低了视觉质量和几何一致性。

与上述方法聚焦~TDOM~生成不同，Gaussian Building Mesh（GBM）（2025.01）~\cite{2501_00625} 将目标从正射影像产物扩展到完整的~3D~建筑网格提取——实现从名称或地址自动提取任意建筑~3D~网格的端到端管线，将~3D~建筑重建从需要专业~UAV~数据采集的专家任务转变为仅需建筑地址的自动化服务。GBM~的研究动机源于~GS2Mesh~在遥感建筑提取场景中的系统性失效：GS2Mesh~的掩码生成模块和预训练立体深度估计模块在遥感视角下经常产生不准确的分割和深度，这促使作者对~2DGS~和分割流程进行针对性改进。

GBM~的管线可以按四个组件读。
第一是 Google Earth Studio~数据采集：输入目标建筑的地址、邮编、名称或经纬度坐标，利用~Google Earth~的全球~2D/3D~数据生成环绕目标的圆形多视角图像序列。
这一步的增量是把现场~LiDAR~或相机采集替换为公开地图数据。

第二是改进的~2DGS~重建。
GBM~选择~2DGS~而非~3DGS~作为骨干，是因为~2D Gaussian Surfel~的平面先验更适合建筑表面。
作者在深度图过滤、平滑处理、训练超参数和网格化超参数上做了工程改造，并通过~masked TSDF integration~直接生成带颜色的~3D~建筑网格。

第三是~GS2Mesh~式网格提取。
系统用训练好的~2DGS~模型为每张输入图像生成立体对，右视图施加水平位移~$[b,0,0]$。
随后，预训练立体匹配网络~DLNR~估计深度图，再经~TSDF~融合和~Marching Cubes~提取网格。
掩码用于过滤单相机可见区域和超范围深度估计，以减少错误面片。

第四是~SAM2+GroundingDINO~前景分割与掩码精炼。
GroundingDINO~根据文本“building”生成边界框，SAM2~利用记忆注意力机制在多视角间传播并精炼掩码。
GBM~再加入形态学操作和~Ramer-Douglas-Peucker~算法，通过多轮重提示优化建筑边界。

在评估方面，GBM~提出基于~360\textdegree~渲染视频的~3D-SSIM~指标——将提取网格的环绕渲染视频与从~Google Earth~训练图像中分割建筑的~360\textdegree~真值视频进行结构相似性比较，提供全序列平均~3D-SSIM~和最小帧~3D-SSIM~两种度量。GBM~还展示了与现成地理注册方法结合后支持建筑信息模型（BIM）、施工验证和自动化安全评估等下游任务的能力。这一管线的核心创新在于端到端自动化——从地址输入到~3D~网格输出无需人工干预，且所有数据源（Google Earth~图像）和工具（2DGS、SAM2、GroundingDINO）均为公开可用。

\begin{takeaway}
文化遗产与测绘将 3DGS 应用于高精度地图生产：正交 splatting 实现真实数字正射影像图生成，增量式 TDOM 管线支持在线更新，2DGS 的平面先验在表面重建中优于 3D 椭球体，端到端管线仅需建筑地址即可自动化提取 3D 网格。
\end{takeaway}

%=====================================================================


\section{其他遥感方法}
\label{sec:20_4}

除卫星~3D~重建、无人机航拍和文化遗产数字化之外，3DGS~在遥感领域还有多种新兴应用，涵盖精准农业、林业调查和生态环境监测等方向。本节从农业与林业遥感、生态环境监测和其他应用~3~个方面展开讨论。

\subsection{农业、林业与行星遥感}
\label{ss:20_5}

本节讨论~2~条技术路线：（1）农业与林业遥感；（2）行星探索与水下遥感。

\subsubsection{农业与林业遥感}

本路线介绍~3DGS~在精准农业和林业调查中的应用，涵盖树木测量、作物表型分析和生物量估计等方法。

3DGS~在精准农业中展现出独特的价值。TreeDGS（2026.01）~\cite{2601_12823} 首次利用~3DGS~作为连续场景表示进行航空树木胸径（DBH，Diameter at Breast Height）测量——传统方法依赖离散点云截面拟合，而~3DGS~的连续密度场允许在任意高度平面进行精确的横截面提取。该方法通过~RaDe-GS~深度感知累积不透明度积分提取密集点集，以不透明度加权实心圆拟合估计~DBH，RMSE~达~4.79 cm，超越~LiDAR~基线（7.66 cm RMSE）。这一突破性结果表明~3DGS~的连续表示在特定几何测量任务中可超越离散传感器直接采集的数据。

与锚点~TreeDGS~聚焦单棵树木的~DBH~测量不同，Wheat3DGS（2025.04）~\cite{2504_06978} 将~3DGS~的测量对象从单一个体扩展到数百个密集排列的麦穗——问题规模从"一棵树的横截面几何"跃迁为"数百个高度相似对象的~3D~实例分割与表型提取"。这一跃迁带来核心困难：麦穗的密集排列、随机分布和外观高度相似性使得跨视角关联极其困难，现有视频追踪器（如~Gaussian Grouping~采用的~DEVA）在稀疏视角和重复结构下完全失效。

为解决这一挑战，本方法提出了"检测-重建-关联-精炼-测量"的自动化流水线：（1）2D~麦穗检测与分割——使用在~Global Wheat Head Dataset~上微调的~YOLOv5~生成每张图像中每个麦穗的边界框，再逐个送入~SAM~获得精确~2D~实例掩码；（2）3DGS~场景重建——以标准~3DGS~重建整个麦田地块的~3D~场景，选择~3DGS~而非~NeRF~正是因为其更直接的场景编辑能力（后续需要对每个高斯基元赋予实例标签）；（3）跨视角实例关联——这是与~TreeDGS~的关键差异所在：TreeDGS~只需在已知位置提取横截面，不存在实例对应问题，而~Wheat3DGS~必须解决"SAM~输出的掩码是实例无关的"这一核心难题。

本方法开发了迭代式匹配-精炼策略：从单视图掩码出发，通过整数线性规划（ILP，Integer Linear Programming）的闭式解优化每个高斯基元的~3D~二值标签$W_k \in \{0,1\}$，将估计的~3D~分割投影到其他视角渲染掩码$\hat{M}'_n$，以精度（precision $> 0.8$）为阈值在~2D~掩码集合中匹配同一物理麦穗（选择精度而非~IoU~是因为单视图分割投影往往存在拉伸，IoU~不可靠），收集多视角掩码后重新求解~ILP~优化~3D~分割，迭代直至所有掩码匹配完毕并为每个高斯赋予唯一的麦穗实例~ID；（4）3D~表型提取——对每个麦穗实例的~3D~点云先经~HDBSCAN~聚类提取主簇和统计离群点去除（SOR，Statistical Outlier Removal）预处理，再从每个麦穗的~3D~点云中提取长度（投影到第一二主成分平面后拟合~2D~平滑样条并计算弧长积分）、宽度（到~1st-2nd-PC~平面的~99th~百分位距离）和体积（Quickhull~算法计算的凸包体积）。

这一方法首次将~3DGS~应用于高通量表型分析，实现了从~2D~图像到~3D~形态测量的全自动管线——从~TreeDGS~证明"3DGS~可精确测量单个对象"到~Wheat3DGS~证明"3DGS~可规模化测量数百个对象"，体现了~3DGS~在农业遥感中从单目标到多目标的扩展。

如果说~TreeDGS~利用~3DGS~的连续密度场在已知树干位置直接提取横截面几何（~DBH），那么~Shen et al.（2024.11）~\cite{2411_08453}~则面对一个不同的任务：油菜生物量估计——这里测量目标不是单棵树的局部截面，而是整株油菜冠层的~3D~体积。这一差异带来了新的技术需求：~TreeDGS~无需分割（树干位置已知），而油菜田间场景包含土壤、标签和相邻植株等多余结构，必须先从~3DGS~重建的场景中分离目标植株才能计算体积。为此，本方法引入了~SAM~辅助的~3D~点云分割管线。数据采集方面，使用~DJI M300 RTK~搭载~Phase One iXM-100~相机（~11664$\times$8750~像素），以相机俯仰角固定$-30^{\circ}$、每~10$^{\circ}$~间隔环绕飞行的方式从~36~个角度采集倾斜影像——这种环绕式多角度策略与~TreeDGS~的航空正射/倾斜互补采集模式形成对比，更适用于平坦农田中低矮冠层的全方位~3D~覆盖。

分割管线的核心是将~SAM~的~2D~分割结果反向投影到~3DGS~点云空间：对~36~视角影像分别以~SAM~生成~2D~分割掩码，将掩码与高斯基元对齐并赋予类别维度，再通过对比学习（contrastive learning）精炼特征一致性，确保不同视角的分割结果在~3D~空间中保持连贯；后处理阶段结合区域生长和噪声滤波去除投影过程中引入的离群点。最终分割精度在苗期和抽薹期合计~293~株上验证，mIoU~达~0.961，F1~达~0.980。体积提取方面，分割后的高斯点云经体素化计算点云体积，再以线性回归拟合地面实测生物量——这一"点云体积模型"显著优于传统基于正射影像的两种基线（地块作物体积模型~$R^2=0.009$、单株作物体积模型~$R^2=0.790$），最终$R^2$~达~0.976，RMSE~为~2.92 g/株，MAPE~为~6.81\%。

在~3D~重建质量方面，3DGS-30k~的~PSNR~达~29.53（~49~分钟训练），接近~Mip-NeRF 360~的~29.72~但训练时间从~45~小时缩短至不到~1~小时，且~FPS~达~91（~Mip-NeRF 360~仅~0.074），兼具精度和实时交互能力。该工作将~3DGS~的农业应用从~TreeDGS~的"利用连续密度场精确测量局部几何"拓展到"重建-分割-体积量化"的完整管线，首次在~296~个油菜品系的大规模田间实验中验证了~3DGS+SAM~组合用于高通量生物量表型分析的可行性。

如果说~TreeDGS~证明了~3DGS~的连续密度场可以精确提取单棵树木的横截面几何，那么~Li et al.（2025.11）~\cite{2511_02207}~则回答了一个不同的问题：能否在重建阶段就将前景植株从背景中分离，直接获得干净的~3D~点云用于表型分析？这一动机源于对当前农业~3D~重建的观察——现有方法（包括~TreeDGS~和~Wheat3DGS）均重建完整场景后通过后处理分割目标对象，背景区域的冗余高斯增加内存消耗、减慢优化速度，并可能通过遮挡或散射效应降低渲染质量；而草莓植株表型分析对重建精度尤为敏感，叶缘锯齿、叶柄弯曲和果实表面等精细结构的保真度直接影响下游形态测量。

为此，本方法提出物体中心（object-centric）3DGS~框架，通过三个协同机制在训练过程中而非训练之后实现前景-背景分离：（1）训练前用~SAM-2~从多视角图像生成~RGBA~输入（Alpha~通道编码前景掩码，区分植株和~10 cm~标定立方体与周围背景）；（2）训练时将~Alpha~掩码作为乘性二值掩码应用于损失计算——给定真值~RGBA~图像$[C, \alpha_{\text{img}}]$和渲染图像$\hat{C}$，损失仅在前景像素上计算$L = \alpha_{\text{img}} \cdot L_{\text{3DGS}}(C, \hat{C})$，$\alpha_{\text{img}} = 0$~的背景像素不产生梯度，对应高斯在自适应精炼阶段（不透明度$\alpha_i < 0.1$时剪枝）被自动移除；（3）每轮迭代使用随机归一化背景颜色$B$合成最终图像$I = C_{\text{splat}} + (1 - \alpha_{\text{acc}})B$（其中$\alpha_{\text{acc}}$为沿射线累积不透明度），进一步抑制模型拟合静态背景。这一物体中心策略与~Wheat3DGS~的~3D~实例分割形成互补——前者在重建过程中隐式分离背景，后者在重建后显式分割实例。

定量对比验证了这一设计的有效性：与后处理背景移除方案（先训练完整场景再用~SAM-2~移除背景）相比，物体中心~3DGS~的~PSNR~提高~4.19 dB、SSIM~提高~0.0132、LPIPS~从~0.0657~降至~0.0550，同时训练速度加快~40\%（6m26s vs. 8m52s），渲染速度提升~2.5$\times$（7.08 FPS vs. 2.83 FPS）——这是因为早期引入背景掩码消除了无关区域梯度，模型将计算完全集中在前景几何上，减少了非信息区域的冗余高斯生成。重建后的点云通过~DBSCAN~聚类分离植株和标定立方体——DBSCAN~无需预定义簇数且能有效处理重建伪影造成的离群点——以立方体的已知尺寸（10 cm）进行尺度标定：对立方体各面拟合多组平行平面，以平均面间距估计重建边长$l_{\text{cube}}$，由此得到全局缩放比$s = 10/l_{\text{cube}}$将~3DGS~坐标转换为物理厘米单位。标定后的植株点云通过~PCA~提取主轴方向，分别计算植株高度（$z$~坐标极差）和冠幅（投影到地面平面后沿主轴的最大横向跨度）。

在~15~棵结果期草莓植株上的表型测量验证表明，植株高度估计的$R^2$~为~0.72、RMSE~为~1.81 cm、MAPE~为~4.85\%，主冠幅（Width~1）的$R^2$~达~0.96、MAPE~为~3.72\%，副冠幅（Width~2）的$R^2$~为~0.90、准确率达~92.94\%——从~TreeDGS~在单棵树木上证明"3DGS~连续表示可超越离散传感器"到本方法在草莓植株上证明"物体中心~3DGS~可在厘米级精度下实现全自动表型提取"，体现了~3DGS~在精准农业中从几何测量到高通量表型分析的深化。

\subsubsection{行星探索与水下遥感}

本路线介绍月球表面~3D~重建、行星探测器重建和水下声纳融合~3DGS~等方法。

3DGS~在极端遥感环境（行星表面和水下）中面临独特的成像挑战：弱纹理、极端光照和长轨迹建图。本节介绍月球表面重建和水下场景重建两个方向的代表性方法。

Casadesus Vila et al.（2026.03）~\cite{2603_18218} 修改了~3DGS~的场景初始化与密度控制环节——用密集深度估计驱动的增量高斯注入替代标准的~SfM~点云初始化和~split/clone~密度控制，同时引入深度监督和语义场损失——标准~3DGS~的可微渲染管线保持不变。

\paragraph{问题}月球表面给~3DGS~带来三重困难：弱纹理导致~SfM~点云极度稀疏，无法提供足够的初始化种子；高对比度光照产生的深阴影在光度上与黑色天空不可区分，标准重建损失~$\mathcal{L}_{\text{recon}}$~无法正确区分二者；更重要的是，标准~3DGS~的全局不透明度重置策略在长轨迹建图中会摧毁已稳定的旧区域高斯，造成灾难性遗忘。

\paragraph{方案}为解决上述问题，该方法提出三项关键修改。（1）前端感知替代~SfM~初始化：用基于~GRU~循环单元的~RAFT-Stereo~网络从立体图像对估计密集深度图，用~CNN~架构的~MANet~对六类月表语义（风化层、岩石、天空、漫游车、人类、着陆器）进行像素级分割；天空像素被显式遮蔽，防止在空区域生成虚假高斯。每帧的深度图和语义标签融合为~3D~语义点云，逐帧增量注入新高斯——每个新高斯的初始尺寸设为与其最近邻点平均距离成正比，自适应覆盖密度不同的区域。（2）密度控制策略重设计：移除标准~3DGS~的全局不透明度重置、split~和~clone~操作（这些操作在长轨迹中不可靠），改用体素过滤策略确保只有未观测区域才添加新高斯，并剪枝与已知~DEM~不一致的高斯。（3）扩展损失函数：在标准重建损失~$\mathcal{L}_{\text{recon}} = (1-\lambda_{\text{SSIM}})\|I-\hat{I}\|_1 + \lambda_{\text{SSIM}}(1-\text{SSIM}(I,\hat{I}))$~和尺度正则化~$\mathcal{L}_{\text{scale}}$~之上，增加四项损失——深度监督损失~$\mathcal{L}_{\text{depth}}$~在有效几何掩码~$M_{\text{dense}}$~上约束渲染深度与估计深度一致：

\begin{equation}
\mathcal{L}_{\text{depth}} = \frac{1}{|M_{\text{dense}}|}\sum_{p \in M_{\text{dense}}} \frac{|D(p)-\hat{D}(p)|}{\max_{q \in M_{\text{dense}}} D(q)}
\end{equation}

即逐像素的归一化~L1~深度误差，分母为当前帧最大深度值以确保尺度不变性。上述损失均通过两个互补掩码施加约束：有效几何掩码~$M_{\text{dense}}$~标识有真实几何的像素（排除天空和超出深度范围的远距点），空区域掩码~$M_{\text{empty}}$~标识确无几何的像素（天空和远距区域）。

在此基础上，空区域正则~$\mathcal{L}_{\text{empty}}$~和密集体积正则~$\mathcal{L}_{\text{dense}}$~协同约束高斯密度分布——前者惩罚空区域中的虚假密度（"天空不应有高斯"），后者鼓励已观测表面充分不透明（"岩石表面应有高斯"）：

\begin{equation}
\mathcal{L}_{\text{empty}} = \frac{1}{|M_{\text{empty}}|}\sum_{p \in M_{\text{empty}}} \alpha_{\text{accum}}(p)^{\beta_{\text{accum}}}, \quad
\mathcal{L}_{\text{dense}} = \frac{1}{|M_{\text{dense}}|}\sum_{p \in M_{\text{dense}}} (1 - \alpha_{\text{accum}}(p))^{\beta_{\text{accum}}}
\end{equation}

其中~$\alpha_{\text{accum}}(p)$~为像素~$p$~处沿射线累积的~alpha~不透明度，超参数~$\beta_{\text{accum}}$~控制正则化强度。直觉上，$\mathcal{L}_{\text{empty}}$~对天空像素处的非零~alpha~施加指数惩罚，$\mathcal{L}_{\text{dense}}$~对几何表面处的低~alpha~施加反向惩罚，二者共同推动高斯聚集到正确位置。

语义~NLL~损失~$\mathcal{L}_{\text{semantic}}$~将语义信息融入~3DGS——每个高斯携带离散语义标签，渲染时按~$\alpha$-混合~累积各类别的概率，再用负对数似然监督：

\begin{equation}
\mathcal{L}_{\text{semantic}} = -\frac{1}{|M_{\text{dense}}|}\sum_{p \in M_{\text{dense}}} \log P_{c(p)}(p)
\end{equation}

其中~$c(p)$~为像素~$p$~的前端分割标签，$P_{c(p)}(p)$~为渲染的该类别概率（计算前~clamp~以保证数值稳定）。

总损失为以上各项的加权求和：

\begin{equation}
\mathcal{L}_{\text{total}} = \mathcal{L}_{\text{recon}} + \lambda_{\text{scale}}\mathcal{L}_{\text{scale}} + \lambda_{\text{depth}}\mathcal{L}_{\text{depth}} + \lambda_{\text{empty}}\mathcal{L}_{\text{empty}} + \lambda_{\text{dense}}\mathcal{L}_{\text{dense}} + \lambda_{\text{semantic}}\mathcal{L}_{\text{semantic}}
\end{equation}

各~$\lambda$~权重系数平衡重建精度、几何约束和语义一致性。

\paragraph{效果}在~LuPNT~模拟器生成的月球南极~120~米遍历路径上，使用单目前向立体相机（$1024 \times 768$，1 Hz），该管线实现~2.8 cm~的均值高度误差，以~5 cm~阈值评估精度~64.3\%、召回~68.6\%，生成约~2400~万高斯（2.4 GB），超越无~LiDAR~的传统点云基线（3.5 cm~误差）。岩石边缘的深阴影区域是最大误差来源（4.0 cm），因为深度估计和语义分割都难以区分阴影岩石边界与黑色天空。3DGS~优化有效平滑了原始点云中的深度噪声和分割边界误差，产出更连贯的场景表示。结果同时表明~MANet~语义分割精度足够高——使用真值语义标签仅带来微小改进，验证了学习型感知模型在月表域的有效性。该方法建立了完整的增量建图管线，为未来集成位姿估计形成完整~SLAM~系统奠定基础。

与锚点~Casadesus Vila et al.~直接修改~3DGS~算法内核（深度驱动的增量高斯注入、语义场损失、体素密度控制）以适应月球表面的极端条件不同，Martínez-Petersen et al.（2026.02）~\cite{2602_13909} 采用了截然不同的策略——不修改~Gaussian Splatting~算法本身，而是构建面向行星机器人探索的完整四阶段自动化管线，核心目标是将现成的~3DGS~变体集成到可复现、可扩展、最小用户干预的标准化重建工作流中。这一设计选择源于不同的研究定位：Casadesus Vila et al.~解决的是"弱纹理、深阴影、长轨迹灾难性遗忘等极端条件下~3DGS~算法如何工作"的问题，而~Martínez-Petersen et al.~解决的是"从原始探测器数据到可交互~3D~视觉模型的端到端工程化落地"问题——前者是算法创新，后者是系统集成。

\paragraph{管线设计}管线的四阶段设计围绕探测器数据的特殊性展开。（1）数据提取与预处理——直接从~ROS2 rosbag~中自动提取~RGB~图像序列、相机内外参和~GNSS/里程计坐标，这一步骤的设计动机在于行星探测器的传感器数据以~ROS~消息格式封装，传统~3DGS~管线无法直接消费。这里一个关键设计是可选的位姿先验注册：当~GNSS~可用时将坐标作为~COLMAP~位姿初始化先验，在~GNSS~拒止场景（如真实行星任务）下则替换为视觉里程计估计——这一灵活性使管线适应有/无绝对定位的两种操作模式。（2）几何重建（Structure-from-Motion）——使用外部~COLMAP~替代~Nerfstudio~默认集成的~SfM~，配置了三种特征匹配策略：序列匹配（仅匹配时序相邻帧，适合有序轨迹但可能遗漏闭环）、穷举匹配（测试所有图像对，覆盖完整但计算昂贵）和混合匹配（结合序列匹配与相机位姿先验引导对应搜索）。在五种轨迹上的对比表明，混合匹配相比~Nerfstudio~默认配置将匹配时间降低~93.1\%、建图时间降低~65.8\%，同时保持全部重建收敛——这一改进的根源在于探测器轨迹的时序有序性使相邻帧特征匹配高度可靠，先验位姿进一步缩小了对应搜索空间。（3）稠密重建与渲染——在~Nerfstudio~中使用~Splatfacto-W~进行~3DGS~建模。选择~Splatfacto-W~而非标准~3DGS~或原始~Splatfacto~的理由基于对探测器数据特征的系统分析：探测器轨迹产生原始~RGB~图像（未经光度标定）、变速运动导致非均匀帧间距、强光照变化使同一区域在不同帧中外观差异显著——这些问题恰好是~Splatfacto-W~从~NeRF-W~继承的鲁棒性机制（外观嵌入、瞬态掩码）所针对的，而标准~3DGS~缺乏这些机制。（4）定量与定性评估——通过重投影误差（平均低于~0.6 px）、每帧平均特征观测数（超过~5000）和平均轨迹长度（超过~7~帧）评估~COLMAP~几何一致性，通过~PSNR（26--29 dB）、SSIM（0.78--0.83）和~LPIPS（0.20--0.26）评估渲染保真度。评估揭示了~Gaussian Splatting~在探测器场景中的固有局限：以直线或近直线为主的轨迹提供不足的视差，导致高斯优化不稳定、对新视角泛化能力差——这一问题与~Casadesus Vila et al.~面临的弱纹理困难形成互补，前者是几何覆盖不足，后者是特征提取困难，但根源相同（行星环境的观测约束）。

\paragraph{工程化特性}整个管线以~Docker~容器化部署，集成~COLMAP、Nerfstudio~和~ROS2~专用环境及全部依赖，支持跨平台和~GPU~集群的无缝迁移。在~ESA BASEPROD~数据集（Bardenas Reales~半沙漠火星类比场地，ExoTer~漫游车采集）的五条轨迹上验证——该数据集并非为摄影测量目的采集，包含运动模糊、非优化视角和不规则帧率，使其成为评估真实探测器约束下重建性能的严格测试平台。实验在配备多块~NVIDIA A6000 GPU、NVMe~高速存储和~10Gbps NAS~的高性能服务器上进行，这一基础设施配置反映了当前行星~3DGS~重建仍需离线处理的现实——论文明确承认管线远超现有行星探测器平台的计算能力，无法实现实时或星载部署。

水下和受限空间遥感方面，Z-Splat（2024.04）~\cite{2404_04687} 与锚点~Casadesus Vila et al.~同样面对传感器基线受限下的~3DGS~几何退化问题，但采取了截然不同的路径——锚点通过密集深度估计和语义场从视觉感知端补偿~SfM~初始化的不足，而~Z-Splat~从~3DGS~的~splatting~投影算子本身出发，引入~z~轴投影以融合声纳的深度轴约束。

\paragraph{问题}锚点~Casadesus Vila et al.~的困难源于弱纹理导致的~SfM~初始化稀疏和长轨迹中的灾难性遗忘——这些都是视觉前端的问题。Z-Splat~关注的是另一个层面的限制：即使在视觉前端正常工作的情况下，当相机基线受限时（水下成像、建筑内部、自主导航等~360\textdegree{}~全景采集不可行的场景），标准~3DGS~的~splatting~操作本身存在结构性盲区。具体而言，标准~3DGS~对每个~3D~高斯执行仿射近似投影，投影后~2D~协方差矩阵为~$\Sigma'_{2D}$（丢弃~z~轴行和列），这意味着沿深度方向的方差~$\sigma_{zz}$~和交叉协方差~$\sigma_{xz}, \sigma_{yz}$~在~splatting~过程中被完全丢弃。从~Fourier-slice theorem~的角度看，多视角相机图像在~Fourier~空间中只采样了若干切片，小基线场景下这些切片无法覆盖深度轴方向的信息，形成~Fourier~空间中的"缺失锥"（missing cone），直接导致~3DGS~沿深度方向的重建质量严重退化。

\paragraph{方案}为填补缺失锥，Z-Splat~在标准~3DGS~的~x-y~平面~splatting~（相机渲染）之外，新增~z~轴方向的~splatting~操作以融合声纳数据。声纳作为主动传感模态发射声脉冲并捕获回波，本质上测量的是场景沿深度轴的密度分布（回波信号可解释为深度的~1D~直方图）。从~Fourier-slice theorem~的角度看，声纳恰好获取了~Fourier~空间中的~z~切片，直接覆盖缺失锥区域。针对两种常用声纳类型，Z-Splat~设计了各自的~splatting~管线：

（1）回声测深仪（Echosounder）：单换能器发射声脉冲，仅能恢复深度信息而无法分辨方位。因此~splatting~将~3D~高斯沿~z~轴投影为~1D~高斯，投影协方差简化为~$\Sigma'_{1D} = \Sigma_{zz}$，仅保留深度方向的均值~$\mu_z$~和方差~$\sigma_{zz}$。渲染时沿~z~轴离散化为深度直方图：对每个深度~bin，将所有可见高斯按~$\alpha$-混合~累积贡献~$Z[i] += T_i \cdot \alpha_i \cdot \exp(-\frac{(z - \mu_z)^2}{2\sigma_{zz}})$，生成与声纳实测回波信号匹配的~1D~深度分布。

（2）前视声纳（Forward Looking Sonar, FLS）：多线阵换能器通过波束成形同时恢复距离和方位角（但不含仰角）——因此~splatting~在~y-z~平面上投影为~2D~高斯，投影协方差保留~$\Sigma'_{2D} = \begin{pmatrix} \sigma_{yy} & \sigma_{yz} \\ \sigma_{yz} & \sigma_{zz} \end{pmatrix}$，相比~echosounder~多了~y~轴维度的信息。由于~RGB~相机的空间分辨率通常高于声纳，y~轴信息主要来自~RGB~相机，因此~FLS~融合的重建精度略优于~echosounder~但不构成质的飞跃——关键增益均在~z~轴方向。

融合通过联合优化实现：相机损失~$L_c = \|I - I_{\text{gs}}\|_1$（L1~图像重建误差）与声纳损失~$L_s = \|S - \hat{S}\|_2$（L2~深度分布匹配误差，echosounder~为向量范数，FLS~为~Frobenius~范数）加权组合~$L = L_c + w \cdot L_s$，权重~$w \in [0.1, 3]$~根据场景纹理丰富程度调节（纹理少的场景增大~$w$~以更依赖声纳）。值得注意的是，与锚点~Casadesus Vila et al.~在前端用~RAFT-Stereo~估计深度后注入~3DGS~的方式不同，Z-Splat~完全不修改~3DGS~的基元初始化或密度控制策略——它仅通过新增~z~轴~splatting~渲染路径和对应的声纳监督信号，在梯度优化层面约束高斯的深度方向参数（$\mu_z, \sigma_{zz}$），标准~3DGS~的可微渲染管线和密度控制保持不变。

\paragraph{效果}通过仿真、模拟和硬件实验验证，RGB-声纳融合相比纯~RGB~的~3DGS~重建显著提升——PSNR~改善~5 dB，Chamfer~距离降低~60\%。所有仿真实验在~NVIDIA 4090 GPU~上平均训练时间约~5~分钟，真实实验中~echosounder~融合~20~分钟、FLS~融合~8~分钟。

与锚点~Casadesus Vila et al.~将语义标签作为月表~3D~建图的辅助监督不同，Qi et al.（2024.12）~\cite{2412_05969} 将语义分割从辅助角色提升为核心目标——在仅有~2\%--6\%标注视角的极端稀疏标签设置下，实现遥感场景的多视角语义分割。这一目标转换带来根本性的设计差异：Casadesus Vila et al.~的语义信息（六类月表标签）由前端~MANet~一次性提供，在优化中作为固定监督信号约束高斯位置；而~Qi et al.~需要解决"少量标注如何驱动全局分割"的问题，语义信息本身成为被优化的对象。

为此，本方法为每个~3D~高斯基元分配~16~维连续语义特征~$s_i$~（替代离散类别值使优化可微），通过与颜色渲染共享的~$\alpha$-混合~管线同时渲染~RGB~图像和语义特征图，再经~MLP~解码器将特征图转换为逐像素类别预测。这种设计使~3DGS~从纯重建工具变为联合重建-分割的统一框架——每个高斯基元同时携带颜色属性和语义属性，空间结构与语义信息在同一个优化过程中相互促进。

然而，稀疏标签设置下（仅~3\%--6\%的视角有标注），标注视图之间的"视觉公共区域"之外存在大量缺乏监督的边界区域。显式点云结构相比隐式神经网络在空间连续性方面更弱——两个相邻高斯基元可能拥有完全不同的语义特征而不受惩罚，导致边界区域分割不一致。本方法从两个层面弥补这一不足：（1）SAM2~伪标签生成——先在某个标注视图上运行~SAM2~生成全分割实例，为每个实例分配伪类别（取该实例与标注标签重叠最多的类别），然后将每个实例掩码作为视觉提示重新输入~SAM2~生成多视角一致的伪标签~$S^p$，以交叉熵损失$L_s^p = \text{CE}(S^{all} \cdot B, S^p \cdot B)$专门约束边界区域~$B$~（~$B=1$~标识边界、$B=0$~标识非边界）；（2）双层聚合损失——在~2D~渲染特征图层面，$L_a^{2D} = \frac{1}{mk}\sum_{i=1}^{m}\sum_{j=1}^{k}s_i \cdot s_j$~鼓励图像空间中相邻特征点的语义相似性（$k=5$~近邻），提升同类区域的特征紧凑度；在~3D~空间点云层面，$L_a^{3D}$~鼓励~3D~坐标中相邻高斯基元的语义特征相似，将~3D~空间连续性注入显式表示。消融实验表明，单独添加~$L_a^{2D}$~提升~mIoU~约~1.1\%，进一步添加~$L_a^{3D}$~再提升~2.4\%，二者互补——前者聚焦~2D~特征紧凑度，后者确保~3D~空间连续性。

整个框架以优化式分割（非前馈推理）运行——30,000~步自监督优化中联合更新颜色属性和语义特征，每~2,000~步执行分裂-合并密度控制同时调整颜色和语义特征的梯度响应。由于伪标签数量远多于真值标签（约~97\%:3\%），真值标签视图与伪标签视图以~1:8~比例采样以平衡监督信号。总损失函数为$L = L_1 + L_{D\text{-}SSIM} + L_s + 0.5 \cdot L_a^{2D} + 0.1 \cdot L_a^{3D}$。在包含~6~个~CARLA~合成子数据集和~3~个~Google Maps~真实子数据集的基准上（每场景约~100~张~$512 \times 512$~图像），本方法在所有子数据集上均达到最高~mIoU——平均超过~IRT（此前最优的隐式方法）~15.5\%，超过~Sem-NeRF~24.42\%，同时优化时间仅约~10~分钟、渲染单视图不足~1~秒，比~IRT~节省~90\%测试时间。这一效率优势源于~3DGS~显式点云结构的一次性渲染策略，避免了~NeRF~类方法逐射线采样和多层~MLP~推理的计算开销。

\begin{takeaway}
其他遥感方法将 3DGS 从纯重建扩展为联合重建-分割的统一框架：连续语义特征和双层聚合损失在仅 3\% 标注视角下实现高质量分割，3DGS 的显式点云结构使渲染效率比 NeRF 方法快一个数量级。
\end{takeaway}

%% ============================================================

\subsection{近期方法脉络：遥感3D生成、遥感二维高斯表示}

这一组近期工作延续 3DGS、2D Gaussian Splatting，主要把改动集中在生成路线、增量改进等环节。下面只保留相对既有路线的变化，便于和前文方法直接对照。

\paragraph{ABot-Earth 0.5}
ABot-Earth 0.5（2026.06）~\cite{2606_09967} 可作为“遥感3D生成”中的新锚点候选。相较于 3DGS，主要改变是：ABot-Earth 0.5：直接以3DGS表示训练生成模型，结合层次LOD和地标增强。3DGS 的显式高斯表示和可微渲染能力保持为基础。
其设计动机是：提出原生3DGS生成框架，学习真实城市重建的几何和纹理分布，在每平方公里10分钟以内生成可Web地图引擎实时可视化的3D环境。
具体而言，模型以地理参考卫星影像为条件，输出3DGS场景表示，并构建分层LOD以支持大范围交互浏览；还探索生成-重建混合的地标增强。
这一锚点判断来自当前分类中的机制关系：候选锚点：生成路线，且相关性 0.95；后续方法可围绕 直接以3DGS作为生成模型输出表示 继续比较和扩展。
适用边界是：生成结果依赖卫星图像信息量；地标细节和几何真实性需要重建增强；评估真实度存在主观性。

\paragraph{GSPan}
GSPan（2026.06）~\cite{2606_17722} 延续 2D Gaussian Splatting 在“遥感二维高斯表示”上的基本设定，主要改变是：GSPan：用2D Gaussian Splatting把各波段残差细节表示为连续可学习Gaussian primitives。3DGS 的显式高斯表示和可微渲染能力保持为基础。
其设计动机是：把固定网格预测改为连续高斯primitive残差表示，实现任意尺度的空间-光谱细节融合。
具体而言，方法设计双流层级交互和空间-光谱注意力，把PAN细节注入各MS波段的2D高斯残差表示。
适用边界是：任务是2D遥感图像融合，不涉及3D场景重建或3DGS渲染。

\section{近期研究（2026 年 8--11 月）}
\paragraph{UAV3DCrop：分开评测作物航拍的外观、几何和冠层高度}
2026.08~\cite{2608_06404}。UAV3DCrop 把重复倾斜航拍的作物重建拆成逐场景优化和零样本前馈两条评测线，分别看新视图、摄影测量参考深度、米制尺度和作物冠层高度。评测保留真实尺度与对齐条件的差异，以免形状相似性掩盖米制误差。

论文发现方法排序随目标而变：Splatfacto-big 的外观表现领先，Scaffold-GS 在深度及冠层高度上更强。当前的新视图划分主要测试插值，深度参考本身依赖摄影测量，在重复纹理、缺少纹理或随风运动的作物上不等同于独立测量真值。

\paragraph{STARS-GS：按结构拆分航拍区域并约束跨区表面}
2026.09~\cite{2609_03447}。STARS-GS 从多视图影像和 SfM 点云出发，按场景结构划区并调整相邻区域边界。相比逐个原语作约束，它把邻近高斯的相对空间组织加入优化，并按照局部几何特征调节表面正则，使分块表面更连续。

GauU-Scene 与 AIR-LONGYAN 用于主要定量和定性比较，UrbanScene3D 与 Mill19 用于补充定性分析。方法研究的是既有航拍数据的表面重建，不直接处理在线跟踪。

\paragraph{SkyAnchor：让无位姿地面序列更新已有米制航拍场景}
2026.09~\cite{2609_13903}。SkyAnchor 将已有米制航空 SfM 和航空高斯作为支架。系统先验证来自地面短序列的候选定位，得到稀疏全局锚点，再在带锚点约束的地面子图中传播尺度和位姿。地面高斯加入时保留航空表示，最后只做有限联合细化，因此不必重新联合求解全部航空与地面相机。

原文在七个有航空和地面采集的场景比较，HorizonGS 场景的连接图像按协议移除。低跨视角重叠区域仍可能产生小的残余错位，通用稠密匹配器的表面覆盖也不能保证米制对齐。结论适用于已经拥有航空几何和高斯支架的更新任务。

\paragraph{HLC-GS：用高度层风险图约束卫星地表模型}
2026.09~\cite{2609_16772}。卫星高斯场在屋顶边界可能同时投影出屋顶和地面，若对高度直接作透明度加权，结果可能落在实际没有表面的中间层。HLC-GS 以高度分散度和主层可靠性生成风险图，只在高风险像素校正不稳的主层，并抑制不一致的远侧或次层响应。

论文在四个 JAX 与三个 IARPA 区域评测有效非森林栅格。相对 EOGS，平均高程 MAE 从 1.46 米降至 1.18 米，RMSE 从 2.78 米降至 2.58 米，2.5 米容差内的比例由 86.09\% 升至 88.61\%。风险图只能约束已有原语，严重遮挡和欠观测下仍难恢复，植被和跨时相变化还可能引入其他误差。

\paragraph{VDGS：按可见性划分航拍场景并补偿增密}
2026.09~\cite{2609_23049}。分块航拍高斯训练中，空间邻近的区域可能有完全不同的相机覆盖。VDGS 从粗模型提取平面原语，结合可见性划分场景，以每个块在图像上的重渲染覆盖决定相机归属。训练时累计低可见锚点的梯度并作补偿，使弱观测位置仍有机会增加原语。

实验覆盖 UrbanScene3D、Mill19 和 MatrixCity 的五个场景。视角不均匀的 Residence 与 Sci-Art 更能体现方法收益，覆盖较均匀的场景中表现主要是有竞争力。论文使用留出相机的图像指标，这与测绘精度是不同的评价目标。

\paragraph{OpenFlyScan：先估计局部重建质量，再补采薄弱区域}
2026.09~\cite{2609_24253}。OpenFlyScan 用图像特征、前馈几何和相机关系预测各区域的高斯重建质量，在正式训练前找薄弱区域并生成补采任务。初始影像和补采影像一起训练高斯，现场反馈由远程工作站计算。

论文在独立 UAVFF3D 实采序列及 Expo West、滴水湖现场数据检查预测和补采效果。该系统不要求无人机配新增硬件，但仍需要地面计算与数据传输。报告的图像质量增益应结合是否增加观测量来比较。

\paragraph{DIBR-GS：用逐步合成的深度伪视图补足遥感观测}
2026.09~\cite{2609_35612}。DIBR-GS 先把单目估计的深度与 SfM 稀疏点对齐，并融合多视图来初始化锚点。表示端加入来自参考图像的跨视角外观特征，训练端逐步把观测深度重投影到未观测视点，为高斯提供附加一致性约束。新的伪视图传递已有几何证据，并未增加真实采集。

作者在 LEVIR-NVS 上比较遥感新视图合成，观察到结构边界和细节的改进。框架仍需逐场景优化，伪视图质量还取决于单目深度和稀疏几何的对齐，不能直接转用为前馈重建。

\paragraph{SemLite-GS：按建筑实例分配航拍高斯训练和网格细化}
2026.11~\cite{2611_19773}。SemLite-GS 将正射影像中的建筑实例身份传播到倾斜视角，用相机和建筑实例的观测关系筛选训练任务。在局部区域并行优化目标高斯时，仍保留固定的周边高斯来解释遮挡。网格阶段将建筑平面约束写入符号距离场细化，尽量保持屋顶和立面结构。

论文使用六个无人机航拍场景，其中五个有 LiDAR 点云用于定量比较。对轻量网格模块单独比较时，固定了实例高斯输入，以隔离网格生成的影响。实验以已配准影像和相机参数为前提，量化结论不包含仅定性评价的第六个场景。


\section*{本章小结}
\addcontentsline{toc}{section}{本章小结}

\begin{futurebox}{本章总结与展望}{ch20_summary}
本章从遥感与航拍角度综述了~3D Gaussian Splatting~的应用方法，涵盖以下核心内容：
\begin{itemize}
\item \textbf{卫星~3D~重建}（\ref{sec:20_1}节）：EOGS~系列框架将~3DGS~适配到~RPC~相机模型和全色数据，ShadowGS~解决多时相阴影不一致性，SkySplat~实现可泛化的~86$\times$~加速重建，From Orbit to Ground~展示了卫星到地面的极端视角外推。
\item \textbf{航空影像重建}（\ref{sec:20_2}节）：端到端~UAV~管线实现了低延迟的实时~3D~重建和~AR/VR~部署，LiDAR~约束和语义分割增强了几何精度和场景理解。
\item \textbf{文化遗产与测绘}（\ref{sec:20_3}节）：CULTURE3D~提供了超大规模基准数据集，3DGS~正交投影生成~TDOM~避免了传统摄影测量管线的复杂性，3DGS~渲染视图上的精确点测量达到厘米级精度。
\item \textbf{其他遥感方法}（\ref{sec:20_4}节）：3DGS~在精准农业（DBH~测量、麦穗表型、生物量估计）、行星探索（月球表面~3D~地图）和水下~3D~重建等新兴方向展现出广阔应用前景。
\end{itemize}

当前遥感~3DGS~面临的主要挑战包括：
\begin{itemize}
\item \textbf{超大规模场景的可扩展性}：城市级重建需要高效的数据分块和分布式优化策略，现有方法在处理数十平方公里区域时仍面临内存和计算瓶颈。
\item \textbf{多时相一致性}：卫星和~UAV~多时相采集中的光照、季节和大气变化对~3DGS~的外观建模提出了更高要求，需要更鲁棒的光照解耦和材质表示。
\item \textbf{几何精度验证}：遥感场景对几何精度（尤其是绝对精度）的要求远高于视觉保真度。3DGS~的~3D~高斯表示在精确表面提取方面仍不如传统~MVS~和~LiDAR~方法。
\item \textbf{实时服务化部署}：从离线优化到在线服务的转变需要解决流式数据摄入、增量更新和边缘部署等工程挑战（参见第~2~章硬件加速和边缘渲染的讨论）。
\end{itemize}
\end{futurebox}
